% !TeX program = xelatex
% 独立源码；无需章节文件。编译两遍以生成目录。
\documentclass[UTF8,a4paper,11pt,oneside,openany,fontset=fandol]{ctexbook}
\usepackage{amsmath,amssymb,amsthm,mathtools}
\usepackage[margin=2.3cm,headheight=16pt]{geometry}
\usepackage{enumitem,booktabs,longtable,array,multicol}
\usepackage[dvipsnames]{xcolor}
\usepackage[most]{tcolorbox}
\usepackage{fancyhdr,listings,tikz}
\usetikzlibrary{arrows.meta,positioning,trees}
\usepackage{hyperref}
\hypersetup{colorlinks=true,linkcolor=MidnightBlue,urlcolor=MidnightBlue,bookmarksnumbered=true,pdfauthor={Federico Prask},pdftitle={博弈论·基础版}}
\definecolor{ink}{HTML}{173A51}
\definecolor{accent}{HTML}{187C82}
\definecolor{soft}{HTML}{F1F6F8}
\definecolor{warn}{HTML}{A46B21}
\ctexset{chapter={format=\Huge\bfseries\color{ink},name={第,章},number=\arabic{chapter}},section={format=\Large\bfseries\color{ink}},subsection={format=\large\bfseries\color{accent}}}
\setlength{\parskip}{0.25em}
\linespread{1.16}
\setlist{nosep,leftmargin=2em,topsep=0.25em}
\allowdisplaybreaks[2]
\emergencystretch=2em
\pagestyle{fancy}
\fancyhf{}
\fancyhead[L]{\small\color{ink} 博弈论}
\fancyhead[R]{\small\nouppercase{\leftmark}}
\fancyfoot[C]{\thepage}
\renewcommand{\headrulewidth}{0.3pt}
\fancypagestyle{plain}{\fancyhf{}\fancyfoot[C]{\thepage}\renewcommand{\headrulewidth}{0pt}}
\theoremstyle{definition}
\newtheorem{definition}{定义}[chapter]
\newtheorem{theorem}[definition]{定理}
\newtheorem{lemma}[definition]{引理}
\newtheorem{corollary}[definition]{推论}
\newtheorem{example}[definition]{例}
\renewcommand{\proofname}{证明}
\newtcolorbox{problem}[1]{enhanced,breakable,colback=soft,colframe=ink!60,boxrule=0.5pt,arc=1mm,title={#1},fonttitle=\bfseries,coltitle=white,colbacktitle=ink,left=2mm,right=2mm,top=1mm,bottom=1mm,before skip=8pt,after skip=6pt}
\newtcolorbox{teach}[1][教学提示]{enhanced,breakable,colback=accent!4,colframe=accent!60,boxrule=0.4pt,arc=2mm,title={#1},fonttitle=\bfseries,colbacktitle=accent,coltitle=white,left=2mm,right=2mm}
\newtcolorbox{pitfall}[1][易错点]{enhanced,breakable,colback=warn!4,colframe=warn!60,boxrule=0.4pt,arc=2mm,title={#1},fonttitle=\bfseries,colbacktitle=warn,coltitle=white,left=2mm,right=2mm}
\newcommand{\Z}{\mathbb Z}
\newcommand{\N}{\mathbb N}
\newcommand{\Pp}{\mathbb P}
\newcommand{\eps}{\varepsilon}
\newcommand{\mex}{\operatorname{mex}}
\newcommand{\SG}{\operatorname{SG}}
\newcommand{\Np}{N^+}
\newcommand{\topbit}{\operatorname{topbit}}
\newcommand{\dep}{\operatorname{dep}}
\newcommand{\Val}{V}
\newcommand{\floor}[1]{\left\lfloor #1\right\rfloor}
\newcommand{\ceil}[1]{\left\lceil #1\right\rceil}
\newcommand{\iva}[1]{\left[#1\right]}
\newcommand{\sol}{\par\noindent\textbf{解析.}\ }
\newcommand{\hint}{\par\noindent\textbf{解题方向.}\ }
\newcommand{\cost}{\par\noindent\textbf{复杂度与实现.}\ }
\newcommand{\src}[1]{\par\noindent{\footnotesize 题目／延伸资料：\url{#1}}\par}
\lstset{language=C++,basicstyle=\ttfamily\footnotesize,keywordstyle=\color{MidnightBlue}\bfseries,commentstyle=\color{gray},stringstyle=\color{accent},showstringspaces=false,breaklines=true,columns=fullflexible,keepspaces=true,frame=single,rulecolor=\color{ink!25},backgroundcolor=\color{soft},tabsize=2}
\begin{document}
\frontmatter
\begin{titlepage}
\setlength{\parindent}{0pt}
\vspace*{1.2cm}
{\large\color{accent}GAME THEORY}\par
\vspace{1cm}
{\fontsize{46}{56}\selectfont\bfseries\color{ink}博弈论}\par
\vspace{0.5cm}
{\Large 精简版}\par
\vspace{0.5cm}
{\color{accent}\rule{\textwidth}{1.2pt}}\par
\vspace{0.3cm}
\begin{minipage}{0.9\textwidth}
\large 麓山国际实验学校 Federico Prask · 2026
\end{minipage}\par
\vfill
{\normalsize }\par
{\small}
\end{titlepage}
\chapter*{阅读说明与课程地图}
\addcontentsline{toc}{chapter}{阅读说明与课程地图}

\begin{pitfall}[适用范围与阅读边界]
本资料面向已经学习基本算法、动态规划、图论与 C++ 的信息学竞赛学生，主题是博弈问题中“最优策略”意义下的胜负判定，第 6 章拓展专题除外（该章引入非公平游戏的双视角算术）。全文默认双方绝对理性、绝不犯错；不覆盖依赖隐藏信息与信念的一般不完全信息博弈，也不讨论随机博弈的胜率优化（仅在纳什均衡一章触及混合策略）。第 6 章为拓展专题（非公平组合游戏与超现实数），属选修提高内容，首次阅读可以整章跳过。

例题采用教学性摘要，不代替评测网站的完整输入输出说明；交互题、提交答案题的工程细节从略。本精简版与完整版的差别在于：各章末尾“分层例题与解析”只保留题面、数据范围与“解题方向”，略去全部“解析”；正文各章的定理与证明保持完整。建议先用精简版独立练习，再对照完整版检验推导。附录的判定模板是教学模板，不是可直接提交的完整程序，使用前应与暴力对拍。
\end{pitfall}

\section*{课程地图}
全书只回答一个问题：\textbf{给定一个博弈，谁赢？}路线是从具体到抽象的四级台阶。

\begin{itemize}
\item \textbf{建立常识}（第 1 章）：把博弈看成一张状态图，胜负从终局倒推；定义 P 态／N 态并给出递推刻画。这是全部结论的地基。
\item \textbf{通用工具}（第 2、3 章）：下模仿棋、策略窃取、设计定式“三板斧”，以及最通用的 PN 态 DP；状态图有环时引入“平局”标签。
\item \textbf{主线武器}（第 4、5 章）：Nim 游戏及其变体是热身，Sprague--Grundy 理论把千姿百态的公平组合游戏折算成一个非负整数（SG 值），和游戏按异或结算。
\item \textbf{认清边界}（第 6、7 章）：撤掉公平性前提，需要 Conway 的超现实数；撤掉“轮流行动”，需要纳什均衡与 Minimax 定理。第 2、3、5、7 章末尾设“分层例题与解析”，第 6 章整章为拓展专题，附录给出可复用的判定模板。
\end{itemize}

\section*{统一记号}
\begin{longtable}{p{0.23\textwidth}p{0.69\textwidth}}
\toprule
记号 & 含义与约定\\\midrule
$A,B$ & 初始状态下的先后手，一旦定下不再变化；具体问题中常用固定的人名。\\
先手／后手 & 当前局面下的下一步行动者／另一方，随局面变化；“把局面做成 P 态”即将先手让给对方。\\
P 态／N 态 & 公平组合游戏中后手必胜态（Previous player wins）／先手必胜态（Next player wins）。\\
$u\to v$，$\Np(u)$ & 状态图的一条合法转移；$u$ 的全部后继构成的集合。\\
$\oplus$ & 按位异或（Nim 和）；$S=a_1\oplus a_2\oplus\cdots\oplus a_n$。\\
$(x)_k$ & 非负整数 $x$ 二进制表示的第 $k$ 位；$\topbit(S)$ 为 $S$ 的最高置位。\\
$\mex(X)$ & 集合 $X$ 中未出现的最小非负整数。\\
$g(u)$，$\SG(u)$ & 状态 $u$ 的 SG 值（Grundy 值），两者完全同义；$g(u)=0$ 恰为 P 态。\\
$G=\{G^L\mid G^R\}$ & 游戏的双视角记号：$G^L,G^R$ 分别为 Left／Right 走一步能到达的局面集合。\\
$\ast$，$\ast_n$ & star：$\ast=\{0\mid 0\}$；一般地 $\ast_n$ 表示 SG 值为 $n$ 的公平游戏（大小为 $n$ 的 Nim 堆）。\\
$\omega$ & 最小无穷序数对应的超现实数 $\{0,1,2,\ldots\mid\}$。\\
$W$ & 双人零和博弈的收益矩阵，$W_{i,j}$ 为玩家 A 的收益；B 的收益为其相反数。\\
纯策略／混合策略 & 固定一种决策／为每种决策分配概率后随机决策。\\
misère（反规则） & “无法行动者胜”的规则族，如反 Nim 游戏；与正常规则结论可能截然不同。\\
$\floor{x},\ceil{x}$ & 下取整与上取整；$\iva{P}$ 为 Iverson 括号，命题成立取 $1$，否则取 $0$。\\
\bottomrule
\end{longtable}

\section*{建议的学习层级}
\begin{longtable}{p{0.12\textwidth}p{0.40\textwidth}p{0.38\textwidth}}
\toprule
层次 & 主线 & 可检查的学习内容\\\midrule
A：基础 & P／N 态递推刻画、PN 态 DP、下模仿棋 & 能对小型状态图手工完成逆向归纳；能说明镜像策略为何永不断裂。\\
B：核心 & 策略窃取、设计定式、Nim 与反 Nim、SG 函数与 SG 定理 & 能独立写出“两段式”证明；能把多堆游戏拆成和游戏并异或结算。\\
C：提高 & Nim-k、广义反 Nim、同步和、阶梯 Nim 应用题 & 能指出错误推广的确切断点；能用整除分块、数位 DP 等工具完成计数。\\
D：专题 & 超现实数、纳什均衡、零和博弈的线性规划与图解法 & 能区分公平／非公平、轮流／同时两套框架；能解 $2\times2$ 零和子博弈。\\
\bottomrule
\end{longtable}

\tableofcontents
\mainmatter

\chapter{博弈问题简介}
\section{一些约定}
在进入具体内容之前，先说明整本书的立场：我们研究的是\textbf{最优博弈}，即假定双方都足够聪明、每一步都采取最优策略、绝不犯错，然后问：从某个初始状态出发，胜负能否被预先确定？这里的“必胜”是严格的数学概念——存在一种策略，使得无论对手如何应对，按此策略行动的一方最终一定获胜。把希望寄托在对手失误上的“策略”不在讨论范围内。

为什么“最优策略”是良定义的呢？关键在于游戏的\textbf{确定性}与\textbf{有限性}：局面向局面的转移完全由玩家的决策决定，且游戏必然在有限步内结束。于是可以把每个合法局面看成一个节点、把一步合法操作看成一条边，得到一张有向图，胜负就能从“终局”倒着递推回初始局面。全篇几乎所有结论，本质上都是在这张状态图上做文章。出于同样的原因，以下约定非常方便：
\begin{itemize}
\item “$A$”“$B$”分别表示初始状态下的先后手，它们一旦定下就不变。
\item “先手”“后手”分别表示\textbf{当前状态}下的先后手：谁下一步行动，谁就是当前先手。推导中常出现“某方把局面做成 P 态”的说法，意思就是把当前先手让给对方，让对方面对一个先手必败的局面。叙述具体问题时通常用固定的 $A$、$B$，而推导胜负时通常用随局面变化的“先手”“后手”。
\item 对于公平组合游戏，“P 态”“N 态”分别表示后手必胜态和先手必胜态（Previous player wins / Next player wins）。P、N 的判定见 1.3 节的递推刻画与第 3 章的 PN 态 DP。
\end{itemize}

\begin{pitfall}[两套称呼不要混用]
“$A$／$B$”是初始的、固定的；“先手／后手”是当前的、随局面翻转的。说“$A$ 把局面做成 P 态”，意思是 $A$ 走完后，轮到行动的 $B$（此时 $B$ 是当前先手）面对必败局面。推导时切换称呼是最常见的笔误来源。
\end{pitfall}

\section{组合游戏}
OI 中研究的博弈问题大部分是\textbf{组合游戏}。这类游戏之所以“好研究”，是因为它的五个特征恰好保证了上一节所说的状态图方法可行、“最优策略”良定义且理论上可以求出。组合游戏的特征：
\begin{itemize}
\item \textbf{完全信息}：任何玩家在所有时候都知道游戏的全部状态。\\
反例：斗地主（玩家看不到其他玩家的手牌）。完全信息使得决策只取决于局面本身，而不取决于“我猜对方有什么”；一旦掺入隐藏信息，策略就与信念、试探和误导挂钩，局面不再能用一张简单的状态图刻画，博弈的难度会陡增（例题 [CF98E] Help Shrek and Donkey 就是一个在少量隐藏信息下仍可求解的例子）。
\item \textbf{无随机性}：游戏进程完全由玩家的决策决定。\\
反例：大富翁（骰子的点数是随机的）。没有随机因素，才谈得上“必胜策略”；一旦允许随机，通常只能讨论“胜率最大的策略”，这属于第 7 章纳什均衡一节的混合策略范畴。
\item \textbf{轮流行动}：玩家严格地轮流进行操作，同一时刻至多一名玩家做决策。这与“所有人同时决策”的静态博弈（见 1.4 节）有本质区别：前者有“状态”和“先后手”，后者没有。
\item \textbf{非合作（零和）}：玩家之间不能合作，通常是两个玩家，一方所得即另一方所失。收益之和为常数使得“我要赢”与“让对手输”完全等价，博弈成为纯粹的对立；一旦出现三人以上或非零和的局面，玩家之间还可能结盟或互相坑害，OI 中几乎不出现。
\item \textbf{确定性的结果}：游戏一定会结束，并根据最终状态产生一个明确的输赢或结果（比如本局得分）。“一定会结束”不是白说的：若允许无限进行下去，就要额外引入“平局”这一第三种结果，判定会复杂得多（例题 [省选联考 2023] 过河卒就是典型例子）。
\end{itemize}

另外要留意“结果”的两种常见形式：一是\textbf{无法行动者负}（全篇绝大多数游戏），二是\textbf{无法行动者胜}（反 Nim 游戏等）。两者只差一个字，结论却可能截然不同；拿到题先看清规则，再决定套用哪套理论。

\section{公平组合游戏}
在组合游戏中，还有一类性质更强的游戏，称为\textbf{公平组合游戏}。

\begin{definition}[公平组合游戏]
一个组合游戏是公平的，若对于同一个局面，无论轮到哪一个玩家移动，其决策空间都是相同的。
\end{definition}

通俗地说，规则不能“认人”：同一局面下甲能走的每一步，乙也一定能走。拿石子、走格子这类游戏天然是公平的；反过来，国际象棋、围棋这类“各自的棋子各自走”的游戏不是公平的——同一个局面下，轮到白方与轮到黑方时能操作的棋子不同。这类游戏双方的地位不对等，需要超现实数这样更精细的工具来刻画（见第 6 章）。

正因为公平，同一局面下两名玩家处于完全对称的位置，于是每个状态只可能属于两类。

\begin{definition}[P 态与 N 态]\label{def:pn}
\textbf{P 态}（后手必胜态）：当前先手无论怎么走，当前后手都有应对策略，使先手最终必败；\\
\textbf{N 态}（先手必胜态）：当前先手存在一步合法走法，使走完后的局面成为 P 态，从而立于不败之地。
\end{definition}

\begin{theorem}[递推刻画]\label{thm:recursion}
在“游戏一定在有限步内结束”的前提下，每个状态恰属于 P、N 之一，且有：没有合法走法的局面是 P 态（当前先手直接判负）；一般地，局面 $u$ 是 N 态当且仅当存在后继 $v\in\Np(u)$ 是 P 态；$u$ 是 P 态当且仅当它的所有后继都是 N 态。
\end{theorem}
\begin{proof}
从终局反向归纳：无后继的状态按定义是 P 态。对一般状态 $u$，若存在 P 态后继，先手走到它即立于不败之地，故 $u$ 是 N 态；若所有后继都是 N 态，则先手无论怎么走都把必胜局面交给对方，故 $u$ 是 P 态。由于每一步转移都严格消耗“剩余步数”，归纳是良基的，每个状态恰好获得一个标签。
\end{proof}

这个递推刻画是第 3 章 PN 态 DP 的理论基础，也是全篇最核心的判定工具，后面几乎所有结论都可以看成它的“手算加速”。

前面已经提到，对于公平组合游戏，状态只分为 P 态和 N 态两类。那么对于非公平组合游戏呢？可以发现状态分为四类：先手必胜、后手必胜、玩家甲必胜、玩家乙必胜——因为“对谁有利”取决于轮到谁走，还取决于“棋子的归属”，两者独立组合出四种结果。

对于感兴趣的选手，超现实数是研究非公平组合游戏的一个有力工具。接下来要讨论的 Sprague--Grundy 理论正是它在“公平游戏”这个特殊情况下的投影；完整的图景要等第 6 章再展开。

\section{静态博弈}
OI 中研究的博弈问题除了组合游戏外，还有一类是\textbf{静态博弈}问题。
\begin{itemize}
\item \textbf{静态}：所有玩家同时作出决策，最后结算。显然静态博弈问题是无所谓状态的——没有“轮到谁”的先后，也就没有状态图可言。
\end{itemize}
最典型的例子是石头剪刀布：双方同时出拳，同时出完立即结算，不存在“看到对方出拳再应对”。OI 中一些“双方各构造一个答案然后互相对抗”的题目也属于静态博弈。因为缺少先后手结构，静态博弈不能套用 P／N 态，需要的是第 7 章“纳什均衡”那一套工具。

每个静态博弈问题都能表示成一个超长方体。具体地，考虑一个 $n$ 名玩家的静态博弈，第 $i$ 名玩家有 $k_i$ 种决策，则这个静态博弈可以表示成一个 $k_1\times k_2\times\cdots\times k_n$ 的超长方体，超长方体的每个单位立方体里写着一个 $n$ 维向量，表示每名玩家的收益。以最常见的两人博弈为例，它就是一张 $k_1\times k_2$ 的\textbf{收益矩阵}：行对应玩家 $A$ 的决策，列对应玩家 $B$ 的决策，格子 $(i,j)$ 里写两人各自拿到的收益。下文讨论纳什均衡时说的“收益矩阵”就是指这张表。

在静态博弈中，还有一类性质更强的博弈，称为\textbf{匿名静态博弈}。
\begin{definition}[匿名静态博弈]
每名玩家的决策空间相同，且对于任意一种可能的博弈结果，决策相同的玩家的收益相同。
\end{definition}
也就是说，玩家之间除了“自己选了哪个决策”之外没有任何身份区别：收益只取决于“我的决策”和“其余人各自选的决策”，其余人之间可以任意交换身份而不影响我的收益。例如“每人报一个 $1\sim10$ 的整数，报出最众数的人赢”就是匿名的；“走私者与检察官”这类角色分工明确的问题则不是匿名的。

\chapter{充要条件：模仿棋、策略窃取与设计定式}
本节介绍判定“初始状态是 P 态还是 N 态”最朴素的手段：直接构造必胜（或必败）的论证。要证明一个状态是 N 态，只需构造先手的第一步乃至整局策略；要证明它是 P 态，则要说明无论先手怎么走，后手都有办法应对到底。常见套路有三板斧，分别对应下面三个小节：
\begin{itemize}
\item \textbf{下模仿棋}：利用局面的对称性，后手（或先手）“照抄”对手的操作，始终与对手保持某种对称；
\item \textbf{策略窃取}：非构造性地证明“后手不可能有必胜策略”，从而先手必胜；
\item \textbf{设计定式}：预先设计一组对应关系（如匹配、配对），把对手的任意一步都映射成自己的合法回应，并把整个博弈化归到已经会做的问题上（通常是 Nim 游戏）。
\end{itemize}
很多题目表面千差万别，最终的解法思路都会落回这三句话上，下面逐个展开。

\section{下模仿棋}
\begin{problem}{经典问题：保龄球瓶}
地上放着一排 $n$ 个保龄球瓶，一击可以击倒一个瓶子或相邻两个瓶子，不能行动者负。判定初始状态。
\end{problem}

\sol \textbf{初始状态是 N 态。}当 $n$ 是偶数时，$A$ 先击倒正中间的两个瓶子；当 $n$ 是奇数时，$A$ 先击倒正中间的一个瓶子。以后 $A$ 不断操作 $B$ 操作的位置的对称位置即可。

\paragraph{为什么这能保证 $A$ 必胜.} 关键在于第一步之后场面的对称性：剩下的瓶子被分成左右两段，两段长度相等（$n$ 为偶数时各 $n/2$ 个，$n$ 为奇数时各 $(n-1)/2$ 个），关于正中间完全对称。于是 $A$ 采取镜像策略：$B$ 在某一段击倒了某个瓶子或相邻两个瓶子，$A$ 就到另一段的对称位置击倒对称的一组。注意这里的操作“击倒一个或相邻两个瓶子”在复制到对称位置后仍是合法操作，且由于两侧的消耗永远同步，只要 $B$ 能击倒的瓶子还存在，其对称位置的那组瓶子也一定还立着——所以 $A$ 的每一步回应都合法。

\paragraph{严格收尾.} 每步至少击倒一个瓶子，游戏必在有限步内结束；又因为 $A$ 的行动是对 $B$ 的合法行动的回应，$B$ 能走一步，$A$ 就永远能跟着走一步，所以先走投无路的一定是 $B$。用博弈论的话说，$A$ 的“能动次数”不少于 $B$，而失败者是轮到自己却不能动的人，故 $A$ 必胜，初始状态是 N 态。$n=1,2$ 的极端情形也被上述方案覆盖：第一步直接击倒全部瓶子，游戏立即结束。

从本题可以提炼出下模仿棋的适用条件：\textbf{操作对象必须能两两配对，且“对称回应”永远合法}。保龄球问题能镜像，正是因为击倒正中间后，左右两侧的剩余部分互不干扰；若操作会跨越中线（例如可以一次击倒任意一段连续瓶子），镜像策略就会失效。

\section{策略窃取}
\begin{problem}{经典问题：擦因数}
黑板上写着 $1\ldots n$ 的每个数，两人博弈，每人要选择黑板上的一个数，擦掉其所有因数，不能行动者负。判定初始状态。
\end{problem}

\sol \textbf{初始状态是 N 态。}$n$ 很大时我们根本不知道必胜策略的第一步该走什么，甚至无从下手验证。策略窃取的思想是\textbf{非构造性}地证明先手必胜：假设后手存在必胜策略，再说明先手可以“偷”来这个策略为己所用，从而导出矛盾。

先给出通用模板：如果先手能找到一步“走了等于没走”的棋（不改变局面的实质、只消耗自己一个回合），而“后手必胜”的假设成立，那么先手可以先走这步不亏的棋，然后完整照搬假设中的后手必胜策略——这样一来先手必胜，与“后手必胜”矛盾，因此后手根本不可能有必胜策略，初始状态必为 N 态。换言之，“存在一步不亏的棋”就足以证明 N 态，尽管我们仍然不知道真正的必胜策略长什么样。

本题中“选 $1$”就是这步不亏的棋：$1$ 的因数只有 $1$ 自己，选 $1$ 只会把 $1$ 从黑板上擦掉，对集合 $\{2,3,\ldots,n\}$ 没有任何影响。于是分两种情况：
\begin{itemize}
\item 如果 $\{2,3,\ldots,n\}$ 是 P 态，则 $A$ 选择 $1$ 即可——$A$ 相当于白走一步后，把一个 P 态交给 $B$，$B$ 成为当前先手必败。
\item 如果 $\{2,3,\ldots,n\}$ 是 N 态，设该状态下先手要取胜可以选 $x$，则 $A$ 第一手就选 $x$。因为“走 $x$ 能赢”只与走完后的局面有关，与谁走无关，所以 $A$ 抢先走掉 $x$ 后，$B$ 面对的局面恰是“先手已走完必胜第一步的 $\{2,3,\ldots,n\}$”，即一个 P 态，$B$ 必败。这里不必担心 $x$ 的因数包含 $1$：即使 $A$ 选 $x$ 也擦掉了 $1$，局面依然对应“$\{2,3,\ldots,n\}$ 中先手走 $x$ 之后”的同一 P 态，论证不受影响。
\end{itemize}
两种情况都推出 $A$ 必胜，故初始状态是 N 态。注意整个证明自始至终没有给出具体的必胜策略，这正是策略窃取的典型特征：\textbf{只能证“存在”，不能给“方案”}。因此它在 OI 中适合用来推导“判定胜负”类的结论；一旦题目要求输出必胜策略（如例题 [CF1147F] ZigZag Game），通常还需要结合其他构造手段。

\section{设计定式}
\begin{problem}{经典问题：台阶游戏（阶梯 Nim）}
有 $n$ 级台阶，第 $i$ 级台阶上有 $a_i$ 个石子。每次可以选择一级台阶（设为第 $i$ 级台阶），把第 $i$ 级台阶上至少一个石子移到第 $i-1$ 级台阶上（当 $i=1$ 时，移到地面上），不能行动者负。判定初始状态。
\end{problem}

\sol 这个游戏看起来比 Nim 游戏多出“向下搬运”的结构，其实两者只隔一层窗户纸：\textbf{只有奇数级台阶上的石子才真正参与博弈，偶数级台阶只是临时堆放区（垃圾桶）}。下面用“不变量”的语言把两个视角统一起来，设 $S$ 为奇数级台阶石子数的异或和。

\paragraph{后手视角.} 后手 $B$ 的策略是：每走一步都把局面恢复成 $S=0$。按 $A$ 的行动分类：
\begin{itemize}
\item $A$ 动了偶数级台阶 $i$：石子被搬到奇数级 $i-1$，$S$ 被破坏。$B$ 不慌不忙地把同一批石子再往下搬一级（从 $i-1$ 搬到 $i-2$；若 $i=2$ 则直接搬下地面），于是奇数级台阶上的石子分布原样恢复，$S$ 回到 $0$。这一来一回等于 $A$ 往垃圾桶里扔了一趟石子，对真正的局面毫无影响。
\item $A$ 动了奇数级台阶 $i$：石子被搬到偶数级 $i-1$（$i=1$ 时直接落地），这等价于从 Nim 的某一堆取走若干石子，$S$ 变成某个非零值。$B$ 把奇数级台阶看作 Nim 游戏的堆，用 Nim 的最优策略走一步（把一个奇数级台阶上的若干石子搬下去），把 $S$ 重新归零。
\end{itemize}
两类回应都合法且把 $S=0$ 交还给 $A$。归纳可知：只要初始 $S=0$，$B$ 就永远不会无路可走；游戏必然结束（石子总数不断减少），而 $A$ 每次行动后 $B$ 都能行动，因此最后无法行动的一定是 $A$。

\paragraph{先手视角.} 反之，如果 $S\ne0$，先手 $A$ 第一手就按 Nim 游戏的最优策略，把某个奇数级台阶上的石子搬到合适的数量使 $S=0$，然后转用上面“后手”的维持策略。

\begin{theorem}[阶梯 Nim 判定]\label{thm:stairnim}
台阶游戏中，先手必胜当且仅当奇数级台阶石子数的异或和不为 $0$。一句话总结：偶数级台阶是垃圾桶，只看奇数级台阶，胜负判定与 Nim 游戏完全相同。
\end{theorem}

\begin{pitfall}[只有奇数级参与异或]
初学者最常见的错误是把所有台阶（包括偶数级）都参与异或，务必记住只有奇数级有效。例题 [SDOI2019] 移动金币和 Gra 最后都会化归到这个结论上。
\end{pitfall}

\begin{problem}{经典问题：开箱子}
有 $n$ 箱石子，第 $i$ 箱石子有 $a_i$ 个。初始时每个箱子都没开启。有两种操作：
\begin{itemize}
\item 打开至少一个箱子。
\item 选择一个已经打开的箱子，取走至少一个石子。
\end{itemize}
不能行动者负。判定初始状态。
\end{problem}

\sol \paragraph{后手视角.} 先建立两个基本观察。第一，所有已打开的箱子整体构成一个 Nim 游戏：已打开的箱子随时可以取石，未打开的箱子只是“待解锁的储备”，取石规则与 Nim 完全一致。第二，后手 $B$ 给自己设定的不变量是：自己每步行动结束后，所有已打开箱子的石子数的异或和为 $0$。

在“不存在异或和为 $0$ 的非空子集”的前提下，$A$ 无论何时打开一批新箱子，这一整批的异或和都不为 $0$（否则这批箱子本身就是一个异或和为 $0$ 的非空子集，矛盾）。于是 $B$ 的回应总是可行的，按 $A$ 的行动分两类：
\begin{itemize}
\item $A$ 打开了一批新箱子：设新批的异或和为 $x\ne0$，则开箱后全体已打开箱子的异或和变成 $x$（原来为 $0$）。$B$ 把已打开箱子当作 Nim 局面，用 Nim 最优策略取石，把异或和归回 $0$。
\item $A$ 从某个已打开的箱子取石：这一步只改变一堆，异或和必变为非零（理由见 Nim 游戏的证明 2，即定理~\ref{thm:bouton} 的证明），$B$ 同样一步取石归零。
\end{itemize}
注意 $B$ 的归零操作只动一个已打开箱子，必然合法；且上面第一条用到条件“任何非空子集异或和不为 $0$”——这正是本题 P 态判定的核心。用归纳可以确认：$B$ 的每一步都有回应，并把 $S=0$ 的局面交还给 $A$。还要排除一种意外：$A$ 会不会一步把局面清空从而取胜？$A$ 一步只动一个已打开的箱子，若 $B$ 归零后的局面只剩一个非空已开箱子，其异或和就等于该箱石子数、不可能为 $0$，所以这种情况不会出现；$A$ 永远无法一步清场，而石子总数有限，游戏最终在 $B$ 的某次回应后结束（全部箱子打开且取空），此时轮到 $A$ 无路可走。

原文把 $A$ 历次开箱称为“第一批”“第二批”，意思正是：无论 $A$ 开多少批，每一批之后 $B$ 都能把不变量补回来。

\paragraph{先手视角.} 反之，如果存在 $a$ 的异或和为 $0$ 的非空子集，则 $A$ 开局直接打开一个极大的这样的子集（不能再并入任何未开箱子而保持异或和为 $0$），把已开箱子的异或和做成 $0$。此后 $B$ 只要打开新的一批 $U$，由极大性知 $U$ 的异或和必不为 $0$（否则 $S^*\cup U$ 是更大的异或和为 $0$ 的子集），$A$ 就按 Nim 策略归零；$B$ 若取石，$A$ 同样归零。角色与后手视角完全互换，于是 $A$ 必胜。

\begin{theorem}[开箱子判定]\label{thm:boxnim}
初始状态是 P 态，当且仅当 $a$ 的任何非空子集的异或和都不为 $0$（即 $\{a_i\}$ 在 $\Z_2$ 上线性无关）。
\end{theorem}

本题与台阶游戏的思路一脉相承：都是先手的“开启动作”会打破不变量，而后手（或赢得主动权的一方）负责把不变量修回来。区别只在于这里的“开启”不是搬运而是解锁，因此判定的条件从“奇数级异或和”变成“是否存在异或和为 $0$ 的子集”。

\section{分层例题与解析}

以下例题与本章的三板斧对应，两题的解析最终都落在“设计定式”上；解析默认读者已读过本章相应小节。

\begin{problem}{[CF1147F] ZigZag Game}
这是一道交互题。

给定一个包含 $2n$ 个节点的完全二分图，每一侧各有 $n$ 个节点。节点 $1$ 到 $n$ 位于二分图的一侧，节点 $n+1$ 到 $2n$ 位于另一侧。你还会得到一个 $n\times n$ 的矩阵 $a$，描述了边的权值。$a_{i,j}$ 表示节点 $i$ 与节点 $j+n$ 之间的边的权值。每条边的权值都不相同。

Alice 和 Bob 在这个图上玩一个游戏。首先，Alice 选择自己要扮演“递增”还是“递减”，Bob 获得剩下的选择。然后，Alice 可以把一个棋子放在图上的任意一个节点。接着，Bob 沿着该节点的任意一条相邻边移动棋子。之后，双方轮流进行如下操作，Alice 先手。

当前玩家必须将棋子从当前顶点移动到某个相邻且未访问过的顶点。设 $w$ 为上一次经过的边的权值。如果玩家扮演的是“递增”，则所经过的边的权值必须严格大于 $w$；如果扮演的是“递减”，则必须严格小于 $w$。无法行动的玩家判负。

给定 $n$ 以及图中各边的权值。你可以选择扮演 Alice 或 Bob，并将与评测机对战。你必须赢得所有对局，答案才会被判为正确。

数据范围：$n\le50$。
\end{problem}

\hint 先想清楚该扮演谁：评测机的先手选项更多，稳妥的是扮演后手；再把“永远有路可走”落实成一组可以预处理出来的完美匹配，匹配的性质与稳定婚姻完全一致。

\cost 稳定性与贪心构造都是 $O(n^2)$ 的（$n\le50$ 绰绰有余），实现时把每条边的权值看成“左部点对右部点的排序”即可。

\src{https://codeforces.com/problemset/problem/1147/F}

\begin{problem}{[Luogu P8851] T2nz.}
这是一道交互题。

小 X 陷入了一个奇怪的梦。在梦境里，她在和小 Q 下一种奇怪的棋。

这是一个 $2^{2n}\times 2^{n}$ 的棋盘，小 X 执黑先行，小 Q 执白后行。

每次操作，需要在当前未满的第一行内，任意选择一格下棋。一格内只能有一个棋子。

下满之后，共有 $2^{2n}$ 行棋子，小 X 的得分为本质不同的行数。

小 X 想最大化她的得分，但小 Q 想最小化小 X 的得分。

你的任务是，扮演小 X 或小 Q，最大化或最小化得分。

若你是小 X，在满足最大化得分 $ans$ 的同时，你也要最大化前 $ans$ 行中本质不同的行数。

数据范围：$n\le7$。
\end{problem}

\hint 每行恰好有 $n$ 个黑子与 $n$ 个白子，一行的样式由黑白对的选择决定；后手用完美匹配压上限，先手用“每步排除一半候选行”的抽屉原理抬下限。

\src{https://www.luogu.com.cn/problem/P8851}

\chapter{PN 态 DP}
把每个状态标成 P 态或 N 态的动态规划，称为 \textbf{PN 态 DP}。它是公平组合游戏题最通用的工具：只要状态和转移可以完整枚举，胜负判定就能机械地算出来；前面几节那些漂亮的定式，可以理解为 PN 态 DP 在特殊结构下的“手算加速版”。

\section{从递推刻画到动态规划}
回顾定理~\ref{thm:recursion} 给出的递推刻画。设 $f(u)=1$ 表示状态 $u$ 是 N 态，则：
\begin{itemize}
\item 若 $u$ 没有合法后继，当前玩家直接判负（无法行动者负的规则下），$f(u)=0$；
\item 一般地，$f(u)=1$ 当且仅当存在后继 $v$ 使 $f(v)=0$（先手走到 P 态即必胜）；
\item 等价地，$f(u)=0$ 当且仅当所有后继 $v$ 都有 $f(v)=1$。
\end{itemize}
若状态图是 DAG，按拓扑序（或者对每个状态记忆化搜索）从后往前递推即可，复杂度为 $O(\text{状态数}+\text{转移数})$。常见的题目里每个状态的出度只有 $O(1)$，状态数在 $10^6\sim10^7$ 级别时直接 DP 就能通过。

\begin{center}
\begin{tikzpicture}[every node/.style={draw=ink!50,rounded corners,fill=soft,font=\small,inner sep=5pt},>=Stealth]
\node (d) at (0,0) {$d$（P）};
\node (a) at (3,1.1) {$a$（N）};
\node (b) at (3,-1.1) {$b$（N）};
\node (c) at (6,0) {$c$（P）};
\draw[->] (a)--(d);
\draw[->] (b)--(d);
\draw[->] (c)--(a);
\draw[->] (c)--(b);
\end{tikzpicture}\\[2pt]
{\small 状态图上的逆向归纳：终局 $d$ 无后继，为 P 态；$a,b$ 有 P 态后继，为 N 态；$c$ 的全部后继都是 N 态，故为 P 态。}
\end{center}

因此这类题目的难点几乎都不在 DP 本身，而在两个方面：
\begin{itemize}
\item \textbf{压缩状态}：必须抓住博弈的本质，把显然等价的局面合并。例如例题“长条格子”把连续几万步的操作压缩成 $O(n)$ 个状态，“FollowingNim 改”把整堆的博弈预处理好再参与 Nim 判定，都是状态设计的典范。
\item \textbf{处理环}：状态图一旦有环，上述递推就不封闭——环上的状态既推不出 P 也推不出 N，需要引入第三种结果“平局”并用图论方法求解，见例题 [省选联考 2023] 过河卒。
\end{itemize}

另外，很多题目在 DP 之前需要先用“策略窃取”式的论证删掉对手的部分选项，把看似庞大的状态空间压到可枚举的规模。PN 态 DP 的扩展性很强，本章末的分层例题（过河卒、长条格子、FollowingNim 改等）会反复展示这些套路。

\section{分层例题与解析}

以下例题的共同骨架是状态图：先想清楚状态与转移，再决定 DP、拓扑还是队列。

\begin{problem}{[省选联考 2023] 过河卒}
题意摘要：$n\times m$ 棋盘上有三枚棋子，双方轮流按规则移动其中一枚，率先无路可走者负；游戏可能无限进行，需要判定每个初始局面的胜负或平局。完整题面见原题（洛谷 P9169）。
\end{problem}

\hint 状态是三枚棋子的位置，图不是 DAG；先在反图上拓扑排序定“胜／负／未定”，再用调整法说明未定点全是平局，最后按步数从小到大用队列求出分出胜负的步数。

\cost 建图与两轮拓扑排序都是 $O(\text{状态数}+\text{转移数})=O(n^3m^3)$（状态数约 $1.6\times10^{6}$ 级别，可过）。本题是“带环状态图的 PN 判定”的教科书例题，框架值得背下来：先拓扑定胜负平，再队列算步数。

\src{https://www.luogu.com.cn/problem/P9169}

\begin{problem}{长条格子}
有一个 $1\sim n$ 的长条格子，第 $i$ 个格子有正整数 $a_i$。假设一个 bot 初始在 $s$，$A$ 和 $B$ 轮流操作，$A$ 先手，每次可以进行其中一个行动（假设 bot 当前位置为 $p$）：
\begin{itemize}
\item 将 bot 移动到 $[p+1,\min(p+k,t)]$ 中的任意一个格子；
\item 给 $a_p$ 减 $1$（只有 $a_p>0$ 时才能进行此行动）。
\end{itemize}
其中 $k$ 是一个给定的参数，$t$ 是 bot 的终点。无法行动的人会输掉游戏。

给出了 $q$ 次操作，每次操作是以下两种之一：
\begin{itemize}
\item \texttt{1 l r d} 表示对于所有正整数 $i\in[l,r]$，将 $a_i$ 变为 $a_i+d$。
\item \texttt{2 s t} 表示当 bot 初始点在 $s$，终点在 $t$ 时，谁有必胜策略。
\end{itemize}
数据范围：$1\le n,k,q\le2\times10^{5}$，$1\le l\le r\le n$，$1\le s\le t\le n$，$0\le a_i\le10^{4}$，$0\le d\le300$。
\end{problem}

\hint 先用镜像论证说明“减 $1$”操作不参与状态、只贡献奇偶条件；再观察 PN 态 DP 中 false 点构成的稀疏链，从右往左跳跃定位，用分块把单次查询压到 $O(\sqrt n)$。

\cost 这个过程可以通过分块优化：把“从 $p-k-1$ 向左第一个偶数”预处理成一步跳跃（或并查集式跳表），每步跳跃至少前进 $k+1$，再用分块把连续跳跃压缩，使每次查询降至 $O(\sqrt n)$。时间复杂度 $O(q\sqrt n)$，配合区间加的维护（块内打标记即可处理 $d\le300$ 的修改）即可通过 $2\times10^{5}$ 的数据。

\begin{problem}{FollowingNim 改}
有 $n$ 堆石子，第 $i$ 堆大小为 $a_i$，每次可以选一堆石子取走至少一个，不能行动者负。

有一个额外限制：如果上一名玩家取走的石子数量 $x\in S$，则下一名玩家必须选择与上一名玩家同一堆石子。

数据范围：$a_i\le10^{4}$。
\end{problem}

\hint 新元素是“锁定”：取到 $S$ 中的数量就把双方锁死在同一堆，直到有人取到 $S$ 之外的数解锁。预处理“锁内子局”的胜负表，冷堆压缩成排名后异或判定。

\cost $f,g$ 的预处理是 $O(\max a_i)$ 级别，判定每组堆 $O(n)$，完全满足 $a_i\le10^{4}$ 的限制。

\begin{problem}{[BJOI2018] 双人猜数游戏}
本题为提交答案题。

Alice 和 Bob 是一对非常聪明的人，他们可以算出各种各样游戏的最优策略。现在有个综艺节目《最强大佬》请他们来玩一个游戏。主持人写了三个正整数 $s,m,n$，然后一起告诉 Alice 和 Bob $s\le m\le n$ 以及 $s$ 是多少。（即，$s$ 是接下来要猜的 $m,n$ 的下限。）

之后主持人单独告诉 Alice $m\times n$ 是多少，单独告诉 Bob $m+n$ 是多少。

当然，如果一个人同时知道 $m\times n$ 以及 $m+n$ 的话就能很容易地算出 $m$ 和 $n$ 分别是多少，但现在 Alice 和 Bob 只分别知道其中一个，而且他们只能回答主持人的问题，不能交流。主持人从他们中的一人开始，轮流询问回答者“知不知道 $m$ 和 $n$ 分别是多少”（回答者只能回答知道/不知道）。

为了节目效果，以及显示出 Alice 和 Bob 的绝顶聪明，主持人希望 Alice 和 Bob 一共说了 $t$ 次“不知道”以后两个人都知道 $m$ 和 $n$ 是多少了。现在主持人找到你，希望让帮他构造一组符合条件的 $m$ 和 $n$。

数据范围：$1\le s\le200$，$2\le t\le15$，保证有解。
\end{problem}

\hint 这不是传统意义的“博弈”，而是公共知识推理：每句“我不知道”都在剔除候选对。按时间顺序分层 DP，答案就是“恰好第 $t+1$ 轮被唯一确定”的对。

\chapter{Nim 游戏及其变体}
Sprague--Grundy 理论（习惯上简称 SG 理论）的作用，其实是把 Nim 游戏的结论推广到更多的游戏上。因此我们先研究 Nim 游戏，再研究 Sprague--Grundy 理论。本章的四份“题意”彼此同构：同为取石子，差别只在“每次动几堆”与“无法行动者负还是胜”两处，正好用来观察结论如何随规则变化。

\section{Nim 游戏}
\begin{problem}{Nim 游戏}
有 $n$ 堆石子，第 $i$ 堆有 $a_i$ 个，每次可以选择一堆，取走至少一个石子，不能行动者负。判定初始状态。
\end{problem}

\begin{theorem}[Bouton 定理]\label{thm:bouton}
设 $S=a_1\oplus a_2\oplus\cdots\oplus a_n$。Nim 游戏中先手必胜，当且仅当 $S\ne0$。
\end{theorem}

这个结论是整个 Sprague--Grundy 理论的基石。为什么答案偏偏是异或而不是别的？一个直观的解释是：一次操作恰好改变一堆，而“把一堆从 $a_i$ 减到 $a_i'$”在二进制下等价于翻转若干个二进制位；异或恰好刻画了“每个二进制位上、值为 $1$ 的堆数的奇偶性”如何被一步扰动。至于它为什么能成为判定条件，看下面的两段式证明。

证明分两个部分：N 态的后继中存在至少一个 P 态；P 态的后继全是 N 态。这种“两段式”写法在本书会反复出现（反 Nim、Nim-k、SG 定理同款），本质是定理~\ref{thm:recursion} 递推刻画的直接应用：第一段说 $S\ne0$ 是 N 态，第二段说 $S=0$ 是 P 态。

\begin{proof}[证明 1（$S\ne0$ 的后继中有 P 态）]
设 $k=\topbit(S)$。因为 $S$ 的第 $k$ 位为 $1$，所以存在 $i$ 使得 $(a_i)_k=1$。那么 $a_i\oplus S<a_i$，把第 $i$ 堆石子取到只剩下 $a_i\oplus S$ 个即可。
\end{proof}

\paragraph{白话翻译.} 异或和 $S$ 的最高位 $k$ 为 $1$，说明二进制第 $k$ 位为 $1$ 的堆有奇数个，于是必存在某堆 $a_i$ 的第 $k$ 位也是 $1$。对这堆做“异或替换”$a_i\leftarrow a_i\oplus S$：由于 $S$ 的最高位恰为 $k$，而 $a_i$ 的第 $k$ 位为 $1$，替换后第 $k$ 位变成 $0$ 且更高位不变，故 $a_i\oplus S<a_i$，确实是合法的取石操作；取完后新的异或和为 $S\oplus a_i\oplus(a_i\oplus S)=0$，一步到达 P 态。若只有一堆石子，则 $S=a_1$，$a_1\oplus S=0$，这一手就是把整堆取完。

\begin{proof}[证明 2（P 态的后继全是 N 态）]
操作一次后 $S$ 一定会发生变化。
\end{proof}
\paragraph{补全.} 证明 2 只有一句话，值得写全：一次操作只把某一堆从 $a_i$ 变成 $a_i'<a_i$，其余堆不动，于是新的异或和
\[
 S'=(a_1\oplus\cdots\oplus a_{i-1})\oplus a_i'\oplus(a_{i+1}\oplus\cdots\oplus a_n)=S\oplus a_i\oplus a_i'.
\]
因为 $a_i'\ne a_i$，所以 $a_i\oplus a_i'\ne0$，进而 $S'\ne S$。特别地，从 $S=0$（P 态）出发，任何一步都会把 $S$ 变成非零值，即 P 态的后继全是 N 态。

\paragraph{归纳收尾.} 每步石子总数严格减少，游戏必然结束；先手面对 $S\ne0$ 时按证明 1 走一步，把 $S=0$ 的局面交给后手；后手无论怎么走，$S$ 都会变回非零（证明 2），先手再归零，如此反复。由于石子数有限，最终先手把局面做成全零（$S=0$ 的终局）时后手无路可走。故 $S\ne0\iff$ 初始为 N 态。

\section{反 Nim 游戏}
\begin{problem}{反 Nim 游戏}
有 $n$ 堆石子，第 $i$ 堆有 $a_i$ 个，每次可以选择一堆，取走至少一个石子，\textbf{不能行动者胜}。判定初始状态。
\end{problem}

直觉上，游戏判定胜负的规则被取反，会导致问题本质的变化，所以反 Nim 游戏的结论会与 Nim 游戏截然不同，但仔细分析后可以发现事实并非如此。规则反转真正改变的是终局附近的行为：正常情况下取走最后一颗石子的人获胜，反规则下他却输了；但只要场上还留有一堆超过 $1$ 个的石子，先手就能把“要不要取最后一颗”的主动权握在自己手里。

\begin{theorem}[反 Nim 判定]\label{thm:antinim}
设 $S=a_1\oplus\cdots\oplus a_n$，$c$ 为 $a$ 中大于 $1$ 的数的个数。若 $\max_i a_i\le1$，则先手必胜当且仅当 $S=0$；否则先手必胜当且仅当 $S\ne0$。也就是说，只有 $\max_i a_i\le1$ 的情况下胜负条件被取反，其余情况下胜负条件不变。
\end{theorem}
\begin{proof}
考虑对 $c$ 分类讨论。每一步只改变一堆，所以一次操作后 $c$ 至多变化 $1$，这保证了不同情形之间不会“串台”。
\begin{itemize}
\item $c=0$：所有堆都是 $1$，每步只能整堆取走，游戏恰好进行 $c_0$（总堆数）步，取最后一颗石子的人反而输，故先手胜当且仅当 $c_0$ 为偶数，即 $S=c_0\bmod 2=0$。这正是唯一一处胜负条件被真正取反的地方。
\item $c=1$：考虑策略窃取，先尝试把最大值变成 $1$，不行就变成 $0$。具体地，设唯一的大堆有 $A>1$ 个石子、其余 $1$ 堆有 $t$ 个。把 $A$ 变成 $1$ 后，留给后手的是 $t+1$ 个“$1$ 堆”；把 $A$ 变成 $0$ 后，留给后手的是 $t$ 个“$1$ 堆”。$t$ 与 $t+1$ 一奇一偶，而全“$1$ 堆”的局面先手胜当且仅当堆数为偶数（见 $c=0$），所以两种做法必有一种能把必败局面交给后手。这里的策略窃取很关键，它消除了规则取反带来的影响：先手通过选择“把 $A$ 停在 $1$ 还是 $0$”，强行把游戏导回 $c=0$ 的已解情形。
\item $c\ge2$：无论如何操作，操作后都有 $c\ge1$。因此直接按照 Nim 游戏的策略操作即可：$S\ne0$ 时按 Nim 的证明 1 把某一堆取到 $a_i\oplus S$，把 $S=0$ 交给对手。需要确认这一手不会把局面送成全“$1$ 堆”的陷阱：Nim 的归零手只动一堆，动完后大堆数仍 $\ge c-1\ge1$；即使恰剩一堆大于 $1$（$c$ 变成 $1$），由 $c=1$ 情形的分析，$c=1$ 时 $S$ 不可能为 $0$，而我们现在交出的局面满足 $S=0$，矛盾。所以对手永远接不到“全 $1$ 堆”的便宜局面，只能按正常 Nim 的节奏输掉。
\end{itemize}
配合石子总数递减的终止性，三种情况分别给出了先手的必胜应对（或 $S=0$ 时后手的维持策略），结论得证。
\end{proof}

反 Nim 给我们的教训是：\textbf{规则反转只影响“全是 $1$”这种极端局面}，遇到“无法行动者胜”的题目，先检查有没有让对手没法拿到便宜终局的手段。

\section{Nim-k 游戏}
\begin{problem}{Nim-k 游戏}
有 $n$ 堆石子，第 $i$ 堆有 $a_i$ 个，每次可以选择至少一堆、至多 $k$ 堆，每堆取走至少一个石子，不能行动者负。判定初始状态。
\end{problem}

普通 Nim 每步只许动一堆，Nim-k 放宽到至多 $k$ 堆，可以想见判定条件要从“异或”（模 $2$ 的按位和）升级为“模 $k+1$ 的按位和”——因为每一步在同一个二进制位上，最多能把 $k$ 个 $1$ 同时清零。

\paragraph{第一猜想及其错误.} 考虑推广 Nim 游戏的结论，猜想：设 $S$ 表示 $a_{1\ldots n}$ 的 $k+1$ 进制异或和（把 $a_i$ 写成 $k+1$ 进制后逐位相加、每一位模 $k+1$ 不进位），则先手必胜当且仅当 $S\ne0$。很遗憾，这个结论是错误的，因为一个 P 态可能到达另一个 P 态：比如说 $k=2$，初始有两堆石子，大小分别为 $1$ 和 $2$，此时 $S=0$（$1+2\equiv0\pmod 3$）。如果 $A$ 把两堆石子都取完，则操作后仍然有 $S=0$，矛盾。症结在于：$k+1$ 进制下“一次操作”并不是“只改一位”——把 $2$ 堆整堆取走相当于把它从 $2$ 改成 $0$，在 $k+1$ 进制下低位的 $2$ 直接消失，总和的每一位都被打乱，导致我们连“P 态走不到 P 态”都保证不了。

\paragraph{修复.} 把 $a_{1\ldots n}$ 写成二进制形式，设 $S$ 表示把这些 01 串看成 $k+1$ 进制数后的异或和（即 $S$ 的第 $j$ 位等于“第 $j$ 个二进制位上为 $1$ 的堆数模 $k+1$”），则先手必胜当且仅当 $S\ne0$。\footnote{原稿此处写作“$S=0$”，按上下文（证明 1 是从 $S\ne0$ 调整到 $S=0$）应为 $S\ne0$，整理时已订正。}修复的要点是把按位运算放在二进制位上进行：一次操作至多影响 $k$ 堆，所以每个数位上至多变化 $k$，正好够不到 $k+1$——这就是把模数取成 $k+1$ 的原因。仍用上面的反例验证：$k=2$，两堆 $1,2$，二进制串为 $01$ 与 $10$，逐位模 $3$ 相加得 $S$ 的两位数位为 $1,1$，$S\ne0$，判定为 N 态；事实也确是如此——$A$ 一步取完两堆即可获胜。

\begin{theorem}[Nim-k 判定]\label{thm:nimk}
设 $S$ 的第 $j$ 位等于“$a_1,\ldots,a_n$ 的二进制第 $j$ 位为 $1$ 的堆数模 $k+1$”，则先手必胜当且仅当 $S\ne0$。
\end{theorem}
\begin{proof}
仍分两个部分，这里的 P 态即 $S=0$ 的局面。
\end{proof}

\begin{proof}[证明 1（$S\ne0$ 的后继中有 P 态）]
从高到低考虑 $S$ 的每一位，设当前考虑到了第 $i$ 位，$(S)_i=j$，已经操作过 $cur$ 堆石子，这些堆中石子数的高位都已经严格减小了，因此低位可以任意修改。那么，考虑把 $(S)_i$ 的 $k+1$ 种取值排成一个环，则如果只用 $cur$ 个堆进行调整，可以得到的是包含 $j$ 的一段弧。如果这段弧就包含 $0$，那直接调整即可。否则，设这段弧是一段区间 $[l,r]$（$1\le l\le j\le r\le k$）。考虑先用 $cur$ 个堆把 $(S)_i$ 调整到 $l$。现在，这 $cur$ 堆石子的数量的第 $i$ 位都为 $0$。又因为 $(S)_i=l$，所以至少存在 $l$ 堆石子的数量的第 $i$ 位为 $1$，再操作其中的 $l$ 堆即可。操作后 $cur$ 会变成 $r$，因为 $r\le k$ 所以操作合法。
\end{proof}

\paragraph{拆开读.} 我们的目标是：从高位到低位逐位把 $S$ 的所有数位清零，全程维护“已经动过的堆（$cur$ 堆）的高位已严格变小”。为什么要这个性质？因为一堆只要高位已经变小，其低位就可以自由改动（无论如何改都不影响已处理好的高位），等于获得了“编辑许可”；没动过的堆高位必须是 $0$，否则它的高位早就该被处理。现在看第 $i$ 位：设它当前的数字和为 $j$，$cur$ 堆在第 $i$ 位可以各自任意取 $0/1$，所以能调出的总和覆盖环上含 $j$ 的一段连续弧（长为 $cur+1$）。若弧内含 $0$，直接调整，第 $i$ 位清零，进入下一位。若弧不含 $0$，说明它是一段不环绕的整数区间 $[l,r]$，$l\ge1$：先把 $cur$ 堆的第 $i$ 位全部调成 $0$（总和降到 $l$），此时数字和为 $l$ 意味着“第 $i$ 位为 $1$ 的堆”恰有 $l$ 个（模 $k+1$ 的意义下至少 $l$ 个），而这些堆都不是 $cur$ 中的堆（$cur$ 的堆第 $i$ 位已是 $0$），它们的高位全为 $0$、可以安全改动；于是再操作其中 $l$ 堆，把它们的第 $i$ 位从 $1$ 改成 $0$，第 $i$ 位总和归零。动过的堆数累计为 $cur+l=r\le k$，操作合法，且这 $l$ 堆从此获得“低位自由”，可以继续处理下一位。如此逐位推进，最终得到 $S=0$ 的局面。

\begin{proof}[证明 2（P 态的后继全是 N 态）]
操作一次后 $S$ 一定会发生变化。
\end{proof}
\paragraph{补全.} 一次操作改变 $c$ 堆（$1\le c\le k$）。取这些堆中发生变化的最高的二进制位 $p$：第 $p$ 位以上所有堆都没有变化；而在第 $p$ 位，凡是变化的堆，其这一位必然是从 $1$ 变成 $0$——否则这堆的值不会变小。因此 $S$ 的第 $p$ 位总和变化了 $-c$，而 $1\le c\le k$ 使得 $-c\not\equiv0\pmod{k+1}$，故 $S$ 必变。特别地，$S=0$ 的局面走一步后 $S\ne0$，P 态的后继全是 N 态。这个“看最高变化位”的论证手法也适用于普通 Nim（$k=1$、模 $2$），是理解后面许多证明的钥匙。

\section{反 Nim-k 游戏}
\begin{problem}{反 Nim-k 游戏}
有 $n$ 堆石子，第 $i$ 堆有 $a_i$ 个，每次可以选择至少一堆、至多 $k$ 堆，每堆取走至少一个石子，\textbf{不能行动者胜}。判定初始状态。
\end{problem}

\begin{theorem}[反 Nim-k 判定]\label{thm:antinimk}
采用定理~\ref{thm:nimk} 中的 $S$。若 $\max_i a_i\le1$，则先手必胜当且仅当 $S\ne1$；否则先手必胜当且仅当 $S\ne0$。
\end{theorem}
\begin{proof}[证明概要]
把反 Nim 游戏和 Nim-k 游戏拼起来，思路与反 Nim 完全平行：先数出大于 $1$ 的堆数 $c$，再分两类看。
\begin{itemize}
\item $c=0$（全是 $1$ 堆）：每步至多取走 $k$ 个“$1$ 堆”，等价于两个人轮流从一堆 $c_0$ 个石子中拿走 $1\sim k$ 个、谁拿最后一个谁输的取石子游戏，先手负当且仅当 $c_0\equiv1\pmod{k+1}$。而全 $1$ 局面下 $S$ 的数位只在最低位有一个非零数字 $c_0\bmod(k+1)$，故这正是“$S=1$”，即先手胜当且仅当 $S\ne1$。
\item $c\ge1$：只要场上还有大堆，就按 Nim-k 的正常策略把 $S$ 归零；反 Nim 一节的经验在这里同样适用——“规则反转”只会在终局附近作祟，而只要有大堆存在，取最后一堆的主动权就握在能维持 $S=0$ 的一方手里。
\end{itemize}
\end{proof}

\begin{pitfall}[自造变体先对拍]
反 Nim-k 这个推广结论的正确性并不平凡。Nim-k 的归零操作一次可能同时动多堆，反 Nim 中“一步只动一堆所以大堆数不减”的论证不再成立，归零后局面可能恰好变成全 $1$ 的尴尬形状，因此“拼接”的细节比反 Nim 微妙得多（讲义作者亦未展开完整证明）。竞赛中遇到这类自造变体，稳妥的做法是先写个暴力对拍验证公式，再放心使用；附录的自测程序对有限范围做了验证。
\end{pitfall}

\chapter{Sprague--Grundy 理论}
\section{SG 函数}
SG 函数做的事情，是把任意公平组合游戏“折算”成 Nim 游戏：它给每个状态赋一个非负整数，称为该状态的 \textbf{SG 值}，表示这个状态“相当于”一个多大的 Nim 堆，从而把 Nim 的胜负结论搬到千变万化的游戏上。当然，这种折算是有代价的，可能会损失部分信息。

具体来说，首先，每个公平组合游戏都可以表示成一个 DAG，游戏规则是：初始时某个节点上放置着一枚棋子，每次操作可以把棋子移动到任意一个出点上，不能移动者负。为什么可以这样表示？回顾 1.3 节：公平组合游戏的状态与合法操作天然构成有向图，而“双方决策空间相同”意味着这个图对两名玩家一视同仁，棋子“属于谁”无关紧要。接下来，考虑对每个节点 $u$，把“棋子在 $u$ 上”的游戏状态对应到一个大小为 $\SG(u)$ 的 Nim 堆上——堆的大小就是这个局面的 SG 值，用来衡量它的“战斗值”。

\begin{definition}[SG 函数与 mex]\label{def:sg}
设 $\Np(u)=\{v_1,v_2,\ldots,v_k\}$ 为 $u$ 的后继集合，$\mex(X)$ 表示非负整数集合 $X$ 中未出现的最小非负整数。定义
\[
 g(u)=\mex\{g(v):v\in\Np(u)\},
\]
函数值 $g(u)$ 称为状态 $u$ 的 \textbf{SG 值}（文献中也叫 Grundy 值）。
\end{definition}

先把这一节的来龙去脉说清：整套理论由 Sprague 与 Grundy 在 1930 年代各自独立建立，故得名 Sprague--Grundy 理论；把赋值看成定义在状态集上的函数，就得到定义~\ref{def:sg} 的 SG 函数。注意大小写：Nim 是游戏专名，必须大写（Nim 游戏、Nim 堆、Nim 和），SG 是 Sprague--Grundy 的缩写。做题时看到“SG 值”“Grundy 值”两个词，心里想着同一个东西即可：给公平组合游戏的每个状态赋一个非负整数，值为 $0$ 的状态恰是 P 态；多个独立游戏合成和游戏时，总值为各值之异或（SG 定理，见下）。

\begin{example}[热身：每次取 $1$ 或 $2$ 颗石子]
设大小为 $n$ 的堆的 SG 值为 $f_n$，则 $f_0=0$（无后继），$f_1=\mex\{f_0\}=1$，$f_2=\mex\{f_0,f_1\}=2$，$f_3=\mex\{f_1,f_2\}=0$，$f_4=\mex\{f_2,f_3\}=1$，依此类推周期为 $3$，而“$f_n=0$ 当且仅当 $3\mid n$”恰对应众所周知的结论（$3$ 的倍数时后手必胜）。
\end{example}

$\mex$ 的定义保证了两件事：第一，$0,1,\ldots,g(u)-1$ 这些值都出现在 $u$ 的后继中，即 $u$ 存在对应大小为 $0\ldots(g(u)-1)$ 的 Nim 堆的后继——好比 Nim 中可以把一堆从 $g(u)$ 取到任意更小值；第二，$u$ 不存在对应大小为 $g(u)$ 的 Nim 堆的后继——Nim 中一步不可能把一堆保持原样。这两条恰好复刻了 Nim 游戏的核心性质：每个 SG 值为 $a$ 的状态，其“能到达的后继的 SG 值集合”包含着通往 $0$ 的路（当 $a>0$），却不含 $a$ 自身。

\begin{theorem}[SG 值刻画 P 态]\label{thm:sgzero}
在正常规则（无法操作者输）下，$u$ 是 P 态当且仅当 $g(u)=0$。
\end{theorem}
\begin{proof}[证明（归纳）]
$g(u)=0$ 时后继 SG 值都不为 $0$（全是 N 态）；$g(u)>0$ 时存在 SG 值为 $0$ 的后继（N 态的定义成立）。
\end{proof}

损失的信息在哪里？不能保证 $u$ 不存在对应大小\textbf{大于} $g(u)$ 的 Nim 堆的后继。也就是说，一个“弱”状态也可能有“强”后继，只是这个后继本身是 N 态、理性玩家不会主动走进去而已。在只分胜负的正常规则下这无伤大雅，但一旦规则变化（例如反 Nim 的“不能动者胜”、或一步可以同时动多个子游戏），这些隐藏的后继就可能颠覆结论——这正是后面“广义反 Nim”“广义 Nim-k”两节讨论的麻烦来源。

\begin{pitfall}[SG 值不排除“强后继”]
$mex$ 只管“最小的未出现值”。$g(u)=2$ 的状态完全可以有 $g=7$ 的后继：正常规则下没人会走过去，可一旦规则反转或一步可动多个子游戏，“强后继”就会被利用。判断一个新规则能否沿用 SG 理论，先问一句：证明里哪里用到了“后继的 SG 值都不超过当前值”？其实哪里都没用到——这正是隐患所在。
\end{pitfall}

\section{广义 Nim 游戏}
\begin{problem}{广义 Nim 游戏}
有 $n$ 个独立游戏 $G_{1\ldots n}$，第 $i$ 个的 SG 值为 $a_i$，每次可以选择一个游戏移动一步，不能行动者负。判定初始状态。
\end{problem}

这就是 Nim 游戏的“角色扮演版”：把 Nim 里的每堆石子换成一个任意的公平组合游戏，把“取石子”换成“在这个游戏里走一步”，SG 值扮演堆大小的角色。

\begin{corollary}[广义 Nim 判定]
设 $S=a_1\oplus a_2\oplus\cdots\oplus a_n$。广义 Nim 游戏中先手必胜，当且仅当 $S\ne0$。
\end{corollary}
\begin{proof}[证明思路]
与 Nim 相同的两段式：第一段与 Nim 的证明 1 逐字相同（把“取到 $a_i\oplus S$”理解成“把第 $i$ 个游戏走到 SG 值为 $a_i\oplus S$ 的后继”，而这样的后继由 SG 值的定义保证存在，因为 $a_i\oplus S<a_i$）；第二段因为一步只改变一个游戏、其 SG 值必变，异或和随之改变，同样与 Nim 的证明 2 相同。
\end{proof}

有了广义 Nim 游戏，读者应该已经能猜到下一步：如何证明“$n$ 个游戏之和的 SG 值等于各游戏 SG 值的异或”？这正是下一节“Nim 和”的内容。

\section{Nim 和：SG 定理}
即 Sprague--Grundy 定理，它是组合博弈论的顶梁柱：\textbf{任意有限公平组合游戏的和游戏，其 SG 值等于各游戏 SG 值的异或}。这句话回答了三个问题：为什么 SG 值恰好是个“数”而不是更复杂的对象？为什么异或运算无处不在？以及（对做题人最实际的）为什么可以把大游戏拆成小游戏分别算 SG 值再异或。

先考虑两个独立游戏 $G_1,G_2$ 的和 $G_1+G_2$：玩家轮流选择其中一个游戏并操作一次。$\SG(G_1+G_2)$ 是多少？

\begin{theorem}[SG 定理]\label{thm:sgtheorem}
$\SG(G_1+G_2)=\SG(G_1)\oplus\SG(G_2)$。
\end{theorem}
\begin{proof}
设 $\SG(G_1)=n$，$\SG(G_2)=m$。首先 $G_1+G_2$ 的后继状态中不含 $n\oplus m$，否则不妨设操作了 $G_1$，则 $G_1$ 的后继状态含 $n$，与 SG 值的定义矛盾。

下面证明 $\forall x<n\oplus m$，存在 $G_1+G_2$ 的后继状态为 $x$。即证 $G_1$ 存在后继状态 $m\oplus x$，或 $G_2$ 存在后继状态 $n\oplus x$。即证 $n>m\oplus x$ 或 $m>n\oplus x$。

设 $k=\topbit(x\oplus n\oplus m)$，则 $(x)_k=0$，$(n\oplus m)_k=1$。不妨设 $(n)_k=0$，$(m)_k=1$。分情况讨论：
\begin{itemize}
\item 如果 $\topbit(n\oplus m)>k$：
\begin{itemize}
\item 如果 $n>m$，则 $n>m\oplus x$；
\item 如果 $n<m$，则 $m>n\oplus x$。
\end{itemize}
\item 如果 $\topbit(n\oplus m)=k$：有 $m>n\oplus x$。
\end{itemize}
\end{proof}

\paragraph{直觉压缩.} 第一段是“合游戏的后继 SG 值不含 $n\oplus m$”，因为一步只能动一个游戏，而 SG 值的定义禁止“原地踏步”；第二段是“小于 $n\oplus m$ 的值全都可达”，它依赖二进制的一个基本事实——对任意 $x<n\oplus m$，异或和的最高差异位 $k=\topbit(x\oplus n\oplus m)$ 上 $n,m$ 中必有一方的这一位是 $1$（另一方是 $0$），而这一位为 $1$ 的一方的 SG 值较大，可以把它“调”到需要的值。两段合起来，后继能取到的 SG 值集合覆盖了 $0,1,\ldots,(n\oplus m)-1$ 的全部值、却不含 $n\oplus m$，由 $\mex$ 的定义（只关心最小的未出现值，更大的值即使出现也不影响结论）即得 $\SG(G_1+G_2)=n\oplus m$。归纳到 $n$ 个游戏的和即得广义 Nim 游戏的结论，这也是下一节“验证”的根据。

SG 定理解决了“分治合”三步中的“合”。那么我们只需要考虑如何“分”和“治”：先把大的游戏分成独立的小游戏，再分别解决（对每个小游戏、每个状态求 SG 值）。这个范式是例题 [ABC278G]、[CF1439E]、[HNOI2014] 等题目共同的骨架。

\section{验证广义 Nim 游戏的结论}
考虑 $n$ 个游戏的和：由 SG 定理，和的 SG 值等于各分量 SG 值的异或，即 $S$；由 SG 值的性质（定理~\ref{thm:sgzero} 的“$g=0$ 恰为 P 态”），整个和游戏初始状态是 N 态当且仅当 $S\ne0$。这个结论与广义 Nim 游戏的结论相同——两条路殊途同归，也算是对 SG 定理的一次“体检”。

\section{广义反 Nim 游戏}
\begin{problem}{广义反 Nim 游戏}
有 $n$ 个独立游戏 $G_{1\ldots n}$，第 $i$ 个的 SG 值为 $a_i$，每次可以选择一个游戏移动一步，当 $a_1=a_2=\cdots=a_n=0$ 时当前执步者胜。
\end{problem}

先解释这个定义。若按字面把反 Nim 推广为“无法行动者胜”，那么在和游戏里“无法行动”意味着每个游戏都走到了自己的终点；把“终点态”集体标注成胜利并不优雅，因为一个游戏的终点（SG 值为 $0$）本身是个正常可到达的状态，把它视为“输”会让所有分析都变得反直觉。于是定义改为：当所有分量都归零（$a_1=\cdots=a_n=0$）时，当前执步者胜。这与反 Nim 的“无法行动者胜”一脉相承（在纯 Nim 的情形下，全零恰是无法行动），但它把“轮到自己时全零”变成了赢局而非输局，效果是：谁把局面做成全零，谁就输了——因为紧接着轮到的对手会立刻获胜。

广义反 Nim 游戏的胜负判定条件与反 Nim 游戏的直接推广稍有不同。这样定义的好处是 SG 值为 $0$ 的状态一定是 N 态——在正常规则下“所有游戏都已归零”是输局，这里它却是赢局——但代价是没法再用 SG 定理来验证结论了：SG 定理本身仍然是正确的，但是合并两个游戏不能递归到结果相同的子问题。说得更直白些：SG 定理的递归证明依赖“子游戏的和”与“子游戏本身”之间的同构关系，而这里的胜负规则取决于“是否所有分量同时为零”这一全局性质，把 $G_1+G_2$ 切成两半后，“某一半全零”并不等于“某一半结束”，递归就此断裂。因此本题的验证只能回到反 Nim 的分类讨论思路，而不能指望把每个游戏化简成一个数。

\begin{corollary}[广义反 Nim 判定]
设 $S=a_1\oplus\cdots\oplus a_n$。若 $\max_i a_i\le1$，则先手必胜当且仅当 $S=0$；否则先手必胜当且仅当 $S\ne0$。
\end{corollary}

直觉与反 Nim 如出一辙：$a_i\le1$ 时每个游戏至多还能走一步，胜负退化为“谁先被迫把局面做成全零”；只要某个 $a_i\ge2$，先手就可以把局面控制在异或非零的安全区，把“做成全零”的最后一手留给对方。

\begin{pitfall}[$\max a_i\le1$ 的真实含义]
这里 $\max_i a_i\le1$ 的含义与石子大小无关，而是每个独立游戏“剩余步数的上限”至多一步。SG 值小的游戏也可能还剩很多步可走（回想“强后继”），所以判定的适用对象是“每个分量都临近终点”的局面，不是数值小的局面。
\end{pitfall}

\section{广义 Nim-k 游戏？}
感觉做不了。因为 $u$ 可能存在对应大小大于 $\SG(u)$ 的 Nim 堆的后继，所以对于一个 P 态，操作大于一堆石子后可能到达另一个 P 态。展开说：Nim-k 一步可以同时动至多 $k$ 个游戏，而 SG 值对应关系只能保证“单游戏一步内不出现 SG 值相同的后继”；当多个游戏同时变动时，各自的隐藏后继（那些 SG 值更大的后继）可能互相补偿，使总和不变。反 Nim 一节尚且能靠“一步只动一堆”的性质维持证明，Nim-k 连这个性质都没有，两段式证明的第二段直接失效，故本讲义不收录它的结论——这本身也是一个提醒：\textbf{推广结论前先检查证明的每个关键步骤在推广后是否依然成立}。

\section{小测验：同步和}
\begin{problem}{小测验}
有 $n$ 个独立游戏 $G_{1\ldots n}$，每次对所有可以移动的游戏\textbf{同时}移动一步，不能移动者负。判定初始状态。
\end{problem}

\hint 这题的“和”不再是 SG 意义下的和（那里每次只动一个游戏），而是\textbf{同步和}：每轮所有还没结束的游戏都要被迫走一步。它更像“每局游戏各自按自己的时钟赛跑，谁先在自己的局里无路可走谁输”的并行竞赛，因此胜负不再由 SG 值决定，而由“谁的游戏先耗完”决定。

\sol 观察到对每个 N 态的独立游戏，$A$ 按照必胜策略操作不劣：如果 $A$ 选择走到一个 N 态点，则即使这个游戏结束得最晚，也是 $B$ 赢。

所以，$A$ 的最优策略是：对每个 N 态的独立游戏，按照赢得最慢的必胜策略操作；对每个 P 态的独立游戏，按照输得最快的策略操作。直觉是：N 态的游戏是 $A$ 的“资产”，资产要尽可能长久地续命（把结束时刻拖到其他游戏之后）；P 态的游戏是 $A$ 的“负债”，负债要尽早爆掉（把“轮到 $B$ 时这局已经结束”的局面尽快交给对方）。双方在所有游戏上同时行动，收益互不干扰，所以最优策略就是对每个游戏独立取最优，再比较“哪个游戏最后结束”。

对每个独立游戏的每个状态，DP 预处理出它是 P 态还是 N 态，以及将来还要走多少步。如果步数最长的独立游戏是 N 态则 $A$ 胜，否则 $B$ 胜。请读者自行验证这个结论与“每个游戏独立地按最优策略走到尽头”是一致的。

\section{分层例题与解析}

以下例题的共同骨架是“分、治、合”：识别独立子游戏，对每个子游戏求 SG 值，用 SG 定理异或结算；个别题先判断“能不能合”，不能合时退回直接分类。

\begin{problem}{[ABC278G] Generalized Subtraction Game}
本题为交互题。

给定整数 $N,L,R$。你将与评测系统进行如下游戏：

有 $N$ 张编号为 $1$ 到 $N$ 的卡片放在场上。

先手和后手轮流进行如下操作：选择一组整数 $(x,y)$，满足 $1\le x\le N$，$L\le y\le R$，并且编号为 $x,x+1,\ldots,x+y-1$ 的 $y$ 张卡片都还在场上，然后将这些卡片从场上移除。

不能进行操作的一方判负，另一方获胜。

你需要选择先手或后手。然后，在你选择的回合与评测系统进行游戏，并取得胜利。

数据范围：$N\le2000$。
\end{problem}

\hint 取走一段连续卡片就把局面切成左右两段的和游戏，先写 SG 递推；直接 DP 是 $O(N^3)$，再观察“取走正中间一段后镜像”可以覆盖 $R>L$ 的全部情形，只剩 $L=R$ 需要 $O(N^2)$ 的 DP。

\src{https://atcoder.jp/contests/abc278/tasks/abc278_g}

\begin{problem}{[CF1439E] Cheat and Win}
考虑一个 $[0,10^{9}]\times[0,10^{9}]$ 的网格。如果 $x\ \mathrm{and}\ y=0$，则称格子 $(x,y)$ 是“好格子”。我们以所有好格子为顶点构建一张图，并在所有相邻的好格子之间连边。可以证明，这张图是一棵树。钦定 $(0,0)$ 是根节点。

AB 两人博弈。最初，部分好格子为黑色，其余为白色。每一回合，玩家选择一个黑色好格子及其若干祖先（可以选 $0$ 个祖先），并将这些格子的颜色反转（白变黑，黑变白）。如果某一玩家无法操作（即所有好格子均为白色），则判负。可以证明，这个游戏一定会在有限步内结束。

最初，所有格子都是白色。你会得到 $m$ 对格子，对于每一对格子，将它们之间的简单路径上的所有格子染成黑色。注意，这里是直接染黑，不是反转颜色。

B 想要获胜，于是决定作弊。他可以在游戏开始前多次进行如下操作：选择一个格子，将它到树根的路径上的所有顶点颜色反转。求 B 至少需要做多少次作弊操作。

数据范围：$m\le10^{5}$。
\end{problem}

\hint 四层结构：分形树与深度公式；把“翻转路径”等价改写成“独立棋子”；单枚棋子的 SG 值是 $2^{\dep_u}-1$；把初始 SG 值的二进制里 $1$ 的连续段数修到零。

\cost 时间复杂度 $O(m\log m)$。总结：本题把“博弈”外衣层层剥掉之后，剩下的其实是一个数据结构题——建虚树、树上差分、按深度定位，博弈含量集中在“每枚棋子 SG 值 $=2^{\dep}-1$”这一步推导上。

\src{https://codeforces.com/problemset/problem/1439/E}

\begin{problem}{Gra}
有 $m$ 个格子排成一行，初始时有 $n$ 个格子 $a_1,\ldots,a_n$ 里有棋子。每次可以选择一枚棋子，将它向右移动到第一个没有棋子的格子中，不能移动者负。判定初始状态。
\end{problem}

\hint 把“棋子”翻译成“空格之间的间隔”：空格是墙，墙与墙之间夹着的棋子是台阶上的石子，这是阶梯 Nim 的伪装。

\begin{problem}{[SDOI2019] 移动金币}
一个 $1\times n$ 的棋盘上最初摆放有 $m$ 枚金币。其中每一枚金币占据了一个独立的格子，任意一个格子内最多只有一枚金币。

AB 两人博弈。每次可以选择一枚金币，并将其向左移动任意多格，且至少移动一格。金币不能被移出棋盘，也不能越过其它金币。不能行动者负。

求有多少种初始状态是 N 态。答案对 $10^{9}+9$ 取模。

数据范围：$n\le1.5\times10^{5}$ 且 $m\le50$。
\end{problem}

\hint 先做判定：又是阶梯 Nim，观察对象是金币之间的空隙；判定化为“奇数项空隙异或和是否为 $0$”后，问题变成带异或约束的计数，用数位 DP。

\src{https://www.luogu.com.cn/problem/P5363}

\begin{problem}{[HNOI2014] 江南乐}
小 A 是一个名副其实的狂热的回合制游戏玩家。在获得了许多回合制游戏的世界级奖项之后，小 A 有一天突然想起了他小时候在江南玩过的一个回合制游戏。

游戏的规则是这样的，首先给定一个数 $F$，然后游戏系统会产生 $T$ 组游戏。每一组游戏包含 $N$ 堆石子，小 A 和他的对手轮流操作。每次操作时，操作者先选定一个不小于 $2$ 的正整数 $M$（$M$ 是操作者自行选定的，而且每次操作时可不一样，但是 $M$ 不能大于选中的堆中的石子数量），然后将任意一堆数量不小于 $F$ 的石子分成 $M$ 堆，并且满足这 $M$ 堆石子中石子数最多的一堆至多比石子数最少的一堆多 $1$（即分的尽量平均，事实上按照这样的分石子方法，选定 $M$ 和一堆石子后，它分出来的状态是固定的）。当一个玩家不能操作的时候，也就是当每一堆石子的数量都严格小于 $F$ 时，他就输掉。补充：先手从 $N$ 堆石子中选择一堆数量不小于 $F$ 的石子分成 $M$ 堆后，此时共有 $N+M-1$ 堆石子，接下来小 A 从这 $N+M-1$ 堆石子中选择一堆数量不小于 $F$ 的石子，依此类推。

小 A 从小就是个有风度的男生，他邀请他的对手作为先手。小 A 现在想要知道，面对给定的一组游戏，而且他的对手也和他一样聪明绝顶的话，究竟谁能够获得胜利？

数据范围：$T\le100$，$N\le100$，$F\le10^{5}$，每堆石子数量 $\le10^{5}$。
\end{problem}

\hint “把一堆分成若干堆”是分裂型操作，分裂后各堆独立，是标准的分治合：算单堆 SG 值再异或；转移里的 $M$ 用整除分块压缩，相同 SG 值的偶数个副本互相抵消后每块只剩两个候选。

\cost 时间复杂度 $O(V\sqrt V)$（$V\le10^{5}$），每组游戏的总判定再异或各堆的 $f$ 即可。

\chapter{拓展：非公平组合游戏与超现实数}
\section{公平这个前提用在哪里}
前面的 Sprague--Grundy 理论能成立，靠的是 1.3 节的公平性前提：同一局面下，双方的决策空间相同。这个前提一旦撤掉——国际象棋里白方只能动白子，轮到白方与轮到黑方时能走的棋完全是两本账——整个框架都会塌掉：
\begin{itemize}
\item \textbf{P/N 二分塌了。}1.3 节说过，非公平游戏的局面可能有四种结果：先手必胜、后手必胜、玩家甲必胜、玩家乙必胜。造一个最简单的例子：桌上有三颗石子，两颗只属于甲、一颗只属于乙，轮到谁谁从自己的石子堆里取走任意正数颗（自己的石子取完了就不能动，判负）。甲先手时先取 $1$ 颗留 $1$ 颗，乙只能拿光自己的 $1$ 颗，甲再拿光自己的，乙无路可走；乙先手时乙拿光自己的 $1$ 颗，甲拿光自己的 $2$ 颗，乙依然无路可走。结论是甲必胜，与先后手无关——“你是谁”比“谁先走”更本质，P/N 两个标签根本不够用。
\item \textbf{SG 函数塌了。}“给局面赋一个非负整数”的前提是双方对这个局面拥有同等选择权；一旦棋子认主人，同一个局面在甲、乙眼里是两本不同的账，一个数装不下两个人的视角。
\end{itemize}
所以非公平游戏需要更精细的语言。本节把 Conway 的整套框架补上：它在数学上完整地包含公平游戏这个特例（第 5 章的内容全部可以看作它的退化情形），这也是原稿提到超现实数的原因。路线是：先给游戏一个“双视角记号”，再定义负、和与比较三种运算，最后把“能比大小”的游戏称为数，并看看数如何“做算术”。

\section{双视角记号}
\begin{definition}[游戏的双视角记号]
把一个游戏写成左右两个选项集合的对
\[
 G=\{G^L\mid G^R\},
\]
其中 $G^L$ 是玩家甲（约定称 Left，记 L）走一步能到达的局面集合，$G^R$ 是玩家乙（Right，记 R）的。
\end{definition}
与全篇一致，约定游戏在有限步内结束、无法行动者负。公平游戏恰好是 $G^L=G^R$ 的特例——所以 Nim 游戏与 SG 值全都住在这个记号里。记号本身平凡，不平凡的是在全体游戏上定义三种运算：
\begin{itemize}
\item \textbf{负}：$-G$ 是把 $G$ 的左右选项对调（再逐个取负），含义是“两人角色互换”。$G+(-G)$ 一定是后手必胜：先手在 $G$ 里走什么，后手就在 $-G$ 的镜像位置回应什么，直到先手无路可走——这是“下模仿棋”在记号层面的化身。
\item \textbf{和}：$G+H$ 表示两局并列，轮到谁走、谁挑其中一局走一步；一局结束后只剩另一局。和满足交换律与结合律，配上负运算后“全体游戏”构成一个群：这是后面一切算术的合法性来源，也是全篇“分治合”中“合”的最终形态。
\item \textbf{比较}：$G=H$ 当且仅当 $G-H$ 是后手必胜；$G>H$ 当且仅当 $G-H$ 无论谁先手都是 L 胜；$G<H$ 对称；$G$ 与 $H$ 不可比（记 $G\parallel H$）当且仅当 $G-H$ 是先手必胜。
\end{itemize}
注意“谁先手”由外部局面决定，但“差是后手胜”是游戏的内在性质，与下棋的人无关——这保证了用“值”做算术的客观性。

\section{四个最小的游戏}
从最小的开始看，每个游戏后面的括号里都写着它到底是谁的优势：
\begin{itemize}
\item $0=\{\mid\}$：双方都没有合法操作，先手直接输——零优势；
\item $1=\{0\mid\}$：只有 L 能走一步，且只能走到 $0$——好比桌上有一颗只属于 L 的石子。R 先手时无路可走；L 先手时拿走石子，R 依然无路可走。无论谁先手都是 L 胜，故 $1>0$；
\item $-1=\{\mid 0\}$：把 $1$ 的左右对调，R 必胜，$-1<0$；注意 $-(-1)=1$，而 $1+(-1)$ 见下一条；
\item $1+(-1)=0$：L、R 各有一颗自己的石子。先手拿走自己那颗，后手拿走自己那颗，先手便无路可走——后手胜，恰好对应“$1$ 与 $-1$ 相加抵消”；一般地 $G+(-G)$ 都是后手胜。
\end{itemize}

看起来像在绕圈子，但立刻有收获：$1$ 可以读作“L 领先一步”，$-1$ 是“R 领先一步”，$1+(-1)$ 是“领先被抵消”。把这个直觉推广开：L 独占 $a$ 颗石子（每步可取任意正数颗、只能 L 取）的游戏值就是整数 $a$（把“独占 $a$ 颗”与数 $a$ 相减，差是后手胜，可归纳验证）；若 L、R 各独占 $a$、$b$ 颗，整局等于 $a+(-b)=a-b$，胜负只看差值。

\begin{example}[验证 $a=2,b=1$]
L 先手时先取 $1$ 颗留 $1$ 颗，R 只能拿光自己的 $1$ 颗，L 再拿光自己的，R 无路可走；R 先手时 R 拿光自己的 $1$ 颗，L 接着拿光自己的 $2$ 颗，R 依然无路。两种先后手都是 L 胜，与值 $2-1=1>0$ 一致。
\end{example}

注意这盘棋里没有共享堆、没有异或，博弈内容与 Nim 完全不同，胜负却被“差一个整数”干净地描述——这是“把局面算成数”的第一个成功案例，也是“游戏是一个数”这句话的含义。

\section{四种结果与值的符号}
值的符号与胜负结果的对应可以一次说清：
\begin{itemize}
\item $G>0$（正游戏）：L 必胜，先手后手都不怕；
\item $G<0$（负游戏）：R 必胜；
\item $G=0$（零游戏）：后手必胜；
\item $G\parallel 0$（与 $0$ 模糊）：先手必胜。
\end{itemize}
四行正好把 1.3 节非公平游戏的四种结果（先手必胜、后手必胜、甲必胜、乙必胜）按“值”重新分类。回头看公平游戏的退化：SG 值为 $0$ 的局面是后手胜，即 $G=0$；SG 值非零的局面是先手胜，即 $G\parallel0$——公平游戏从不出现 $G>0$ 或 $G<0$。这从根上解释了为什么公平游戏只需要问“是否为 $0$”：双方的选项完全对称，谁也积累不出“专属优势”，能积累的只有“模糊的先手优势”，于是 Sprague--Grundy 理论只关心一个比特的信息；非公平游戏则不然，L 可以拥有实打实的“正资产”。

\section{数：冷游戏的算术}
游戏世界里存在“不是数”的游戏（下一节的 $\ast$ 就是），Conway 把数定义为如下特殊类型的游戏。

\begin{definition}[数]
左右选项满足“每个左选项都小于每个右选项”的游戏称为\textbf{数}。
\end{definition}

$0=\{\mid\}$ 是第一个数；此后由简到繁依次造出整数（$1=\{0\mid\}$、$2=\{1\mid\}$、$\ldots$）、分数（$\tfrac12=\{0\mid 1\}$：L 的选项是 $0$、R 的选项是 $1$，值恰好卡在两者之间，即“比 $0$ 好、比 $1$ 差”的半目优势）、乃至无穷与无穷小（$\omega=\{0,1,2,\ldots\mid\}$）。全体数构成一个可以加减乘除的全序结构，这就是\textbf{超现实数}。数之间没有模糊性——任意两个数都能比出大小，这与“棋子认主人”造成的混乱正相反，因此取值是数的游戏也叫\textbf{冷游戏}：双方都没有抢着先手的动机，局面的价值像温度计读数一样稳定，先手可能反而吃亏。

冷游戏的好处是可以直接做算术：几个数游戏并成一局，总值就是各数的和；和为正则 L 占优、为负则 R 占优、为零则后手胜。

\begin{example}[$\tfrac12+\tfrac12=1$]
两局 $\{0\mid 1\}$ 并玩：L 先手时把任一局走到 $0$，剩下 $0+\tfrac12=\tfrac12>0$，R 无论怎么走都救不回来；R 先手时把任一局走到 $1$，剩下 $1+\tfrac12=\tfrac32>0$，L 更是稳赢。两个“半目”合成“一目”，算术与博弈严丝合缝。
\end{example}

这就是“超现实数是研究非公平组合游戏的有力工具”的技术内容：把局面翻译成值，把博弈翻译成算术。

\section{模糊游戏：star 与 SG 值}
再看那个“不是数”的游戏：
\[
 \ast=\{0\mid 0\},
\]
双方都只有一步可走，先手走到 $0$ 即胜。$\ast$ 与 $0$ 不可比：任何一方先手都能赢，谁也不能声称自己“更大”——它正是 SG 值为 $1$ 的公平游戏，即大小为 $1$ 的 Nim 堆。更有意思的是 $\ast+\ast=0$：两局 $\ast$ 并玩，先手把其中一局走到 $0$，后手把另一局走到 $0$，先手无路可走——后手胜。对照 SG 值：$1\oplus1=0$，分毫不差。一般地，SG 值为 $n$ 的公平游戏都记作 $\ast_n$（大小为 $n$ 的 Nim 堆），且
\[
 \ast_n+\ast_m=\ast_{(n\oplus m)},
\]
SG 定理（Nim 和）在这套记号下浓缩成一行：公平游戏相加，值按异或合成。异或为什么无处不在，答案就在这里：不是异或神奇，而是公平游戏在“游戏群”里的结构恰好是 $\ast$ 家族的加法——$\ast_n+\ast_n=\ast_{(n\oplus n)}=\ast_0=0$，任何公平游戏与自己相加都是后手胜，与“SG 值异或自己得 $0$”完全一致。

\begin{pitfall}[口径与边界]
这里有一个需要厘清的口径问题。如果把“超现实数”理解成 Conway 的整套“游戏即数”体系，那么 SG 值确实是其中定义清晰的一个大家族——原稿把 SG 值视为超现实数的一个特例，正是取这个广义口径；若取狭义（只指那些能比大小的数），SG 值不是数——$\ast$ 不是 $0$，$\ast_2$ 也不是实数 $2$。“公平游戏”与“数”是这套体系里两个不同的子类：前者左右选项相同、典型地与 $0$ 模糊，后者左小右大、彼此全能排出大小。两种口径都能自洽；关键只是别把两者混为一谈——$\ast$ 的 SG 值是 $1$，但 $\ast$ 本身并不等于数 $1$。

这套框架的边界也要说清楚。它默认正常规则（无法行动者负）与和游戏的封闭性：一旦规则反转（misère 系，如反 Nim 游戏），“差是后手胜即相等”这套公理不再成立，这正是反 Nim 的结论需要另起炉灶、“广义反 Nim”无法再用 SG 定理验证的深层原因；而“一步同时动多个子游戏”之类的变体同样会离开这套框架（见“广义 Nim-k 游戏？”一节的讨论）。另外，对一般的模糊游戏求“值”并不容易：两个游戏比大小在计算上可以非常困难，竞赛题几乎只会碰到冷游戏退化成整数的情形（即独占堆那种例子）。因此本节的实际用途有三：看懂公平游戏为什么只需要 SG 值；遇到“这不是公平游戏”的题目时能立刻认出 SG 不适用；以及在极少数构造题里，用这套记号做小规模的手工推演。
\end{pitfall}

\section{拓展：数从游戏中诞生——单纯性法则}
6.5 节把 $\tfrac12=\{0\mid 1\}$ 当作事实抛了出来：凭什么 $\{0\mid 1\}$ “就等于” $\tfrac12$？本节给出这套造数法的完整规则。它出自 Berlekamp、Conway、Guy 的《Winning Ways for Your Mathematical Plays》第一章（另见 Conway《On Numbers and Games》），是超现实数体系的“施工规范”。

\begin{definition}[出生日]
从 $\{\mid\}$（即 $0$）出发，每用已有的游戏构造一次新的数，就说新数“晚出生一天”：$0$ 出生于第 $0$ 天，$1=\{0\mid\}$ 与 $-1=\{\mid 0\}$ 出生于第 $1$ 天，$\{0\mid 1\}$ 出生于第 $2$ 天，依此类推。
\end{definition}

\begin{theorem}[单纯性法则]\label{thm:simplicity}
设每个左选项与每个右选项都是已构造出的数，且每个左选项严格小于每个右选项，则 $\{L\mid R\}$ 的值是\textbf{严格大于所有左选项、严格小于所有右选项的那个出生最早}的数。
\end{theorem}

严格证明见《Winning Ways》第二章与《On Numbers and Games》；本书把它当“施工规范”使用，重点是算对值。“相等”仍按 6.2 节的定义检验：两个游戏相等，当且仅当它们的差是后手必胜。

\begin{example}[二进分数的诞生]\label{ex:dyadic}
$\{0\mid 1\}$：严格介于 $0$ 与 $1$ 之间、出生最早的数是 $\tfrac12$（第 $2$ 天），故 $\{0\mid 1\}=\tfrac12$，这正是 6.5 节直接验证过的等式。类似地
\[
 \{1\mid 2\}=\tfrac32,\qquad \{0\mid \tfrac12\}=\tfrac14,\qquad \{-1\mid 0\}=-\tfrac12 .
\]
按此规范造出的数先是全体二进分数（分母为 $2$ 的幂），很晚之后才轮到 $\tfrac13$、$\pi$ 乃至无穷与无穷小——超现实数的整个数轴是“从游戏里长出来”的。
\end{example}

\paragraph{伐木游戏（Hackenbush）竹竿.} 单纯性法则最经典的应用对象是\textbf{蓝红伐木游戏}：图由若干条边接在地面（或接在别的边上）构成，Left 只能砍蓝边、Right 只能砍红边；一条边被砍掉时，它上方悬空的整个部分一起倒下；无法行动者负。一根由 $n$ 条蓝边叠成的竖直竹竿，Left 砍最底一条就让 Right 无事可做，值为 $\{n-1\mid\}=n$——这正是 6.3 节“L 独占 $n$ 颗石子”画成图的样子。混色竹竿才是新意。\textbf{蓝底红梢}（自下而上：蓝、红）：Left 砍蓝底，整根倒掉，得 $0$；Right 砍红梢，剩下一条蓝边，值 $1$。于是
\[
 \text{蓝红}=\{0\mid 1\}=\tfrac12 .
\]
再接一条红（自下而上：蓝、红、红）：Left 仍只有砍底一条路，得 $0$；Right 砍中间的红剩值 $1$、砍顶端的红剩 $\tfrac12$，Right 取对自己更小的 $\tfrac12$，故
\[
 \text{蓝红红}=\{0\mid \tfrac12\}=\tfrac14 .
\]
对称地，红底蓝梢的值为 $-\tfrac12$。竹竿每向上多接一节，值的绝对值就折半——二进分数就这样长在竹竿上，6.5 节“两个半目合成一目”的围棋直觉由此获得精确模型。

\section{拓展：开关、热博弈与温度}
若 $x\ge y$ 且 $x,y$ 都是数，则 $\{x\mid y\}$ 不满足“左选项 $<$ 右选项”，因此\textbf{不是数}，称为\textbf{开关}（switch）。开关是“热”的：先手总能从自己的选项里拿到好处。

\begin{example}[开关的两端]\label{ex:switch}
下一节 Domineering 的 $2\times2$ 棋盘值为 $\{1\mid -1\}$：Left 先手走成 $1$，Right 先手走成 $-1$，谁先谁赢、与 $0$ 模糊，记作 $\pm1$。一般地 $\{3\mid -1\}$ 读作“Left 先手可拿到 $3$、Right 先手可拿到 $-1$”的局面。
\end{example}

开关虽然不是数，实践中常用两条粗略刻画（选学内容，不作严格定理）：\textbf{均值} $\tfrac{x+y}{2}$ 大致描述“双方轮流各走一步后平均谁赚”；\textbf{温度} $\tfrac{x-y}{2}$ 大致描述“先手权在这一点上值多少”。由此得到一条实用的启发式——\textbf{先玩最热的分支}：和游戏里有多个分支时，先手应当先把先手权花在温度最高的分支上。

\begin{example}[先玩最热]\label{ex:hotstrat}
$\{3\mid -1\}+\tfrac12$：Left 先手走开关，得 $3+\tfrac12=\tfrac72$，胜；Right 先手也走开关，得 $-1+\tfrac12=-\tfrac12$，胜。双方的最优首着都是开关而不是 $\tfrac12$——先手权花在温度高的地方才不浪费。注意该局面与 $0$ 模糊：均值 $\tfrac{3-1}{2}+ \tfrac12=\tfrac32>0$ 只说明“长期轮流平分下 Left 净赚”，并不决定单局胜负；这正是开关不能当数用的原因。
\end{example}

\begin{center}
\begin{tabular}{cccl}
\toprule
开关 & 均值 & 温度 & 一句话画像\\
\midrule
$\{1\mid -1\}$ & $0$ & $1$ & 纯先手权，谁先谁赢（Domineering $2\times2$）\\
$\{2\mid -2\}$ & $0$ & $2$ & 更烫的纯先手权（兑现支票 $X_2$）\\
$\{3\mid -1\}$ & $1$ & $2$ & 均值偏向 Left 的争夺点\\
\bottomrule
\end{tabular}
\end{center}

再次强调口径：均值与温度是\textbf{启发式坐标}，不是相等关系；两个均值相同的开关未必相等。它们的价值在于给“先玩哪个分支”排序，以及快速估计“这步先手权值几目”。

\section{拓展：微小值 up、down 与 up-star}
数的世界之外还有一群“比 $0$ 精细、又比任何正数都小”的微小值，它们刻画的是先手权以外的第二层优势。

\begin{definition}[up 与 down]
\[
 \uparrow=\{0\mid \ast\},\qquad \downarrow=\{\ast\mid 0\}=-\uparrow .
\]
\end{definition}

以下三件事实都可以用 6.3 节的“分谁先手”手工验证（完整验证见本节例 3）：
\begin{itemize}
\item $\uparrow>0$：无论谁先手都是 Left 胜——$\uparrow$ 是货真价实的正优势，只是很小；
\item $\uparrow\parallel\ast$：微小值与先手权互为模糊，谁也不能压倒谁；
\item $\uparrow+\downarrow=0$：镜像策略的再现——Right 在 $\downarrow$ 中的每一步恰好回应 Left 在 $\uparrow$ 中的那一步。
\end{itemize}

\begin{theorem}[微小性]\label{thm:up}
$\uparrow>0$，且对任意正数 $x$ 都有 $0<\uparrow<x$；$\downarrow$ 对称地满足 $-x<\downarrow<0$。
\end{theorem}
\begin{proof}[证明思路]
对 $x$ 的出生日归纳：正数 $x$ 的左选项中必有非负数，逐一与 $\uparrow$ 的选项比较即可；完整归纳见《Winning Ways》。博弈含义是：$\uparrow$ 是“一丝真实优势”，它永远输给哪怕最小的正数优势，却永远压过 $0$。围棋贴目里“半目”的精确定价用的正是这类微小值。
\end{proof}

\section{拓展：两个经典游戏——Domineering 与蓝红伐木}
\begin{problem}{Domineering（纵横多米诺）}
在方格棋盘上，Left 每轮竖放一张 $1\times 2$ 骨牌，Right 每轮横放一张 $2\times 1$ 骨牌，格子不能重叠，先无法放置者负。求给定棋盘的值。
\end{problem}

它在 1.3 节的意义下显然非公平：同一张空棋盘，轮到 Left 与轮到 Right 时能放的骨牌方向不同。小棋盘可以直接按定义算值：
\begin{itemize}
\item $1\times n$ 单行棋盘：Left 永远无处落子。$1\times2$：Right 放满即空，值 $\{\mid 0\}=-1$；$1\times3$：Right 放完必剩 $1\times1$，右选项为 $0$，值 $\{\mid 0\}=-1$；$1\times4$：Right 的右选项为放中间（断成两个 $1\times1$，值 $0$）与放一端（剩 $1\times2$，值 $-1$），删去被支配的 $0$ 得 $\{\mid -1\}=-2$。归纳可得
\[
 v(1\times n)=-\bigl\lfloor n/2\bigr\rfloor .
\]
\item 单列 $2\times1$：Left 竖放占满即空，Right 无处可放，值 $\{0\mid\}=1$。
\item $2\times2$：Left 竖放剩单列 $2\times1$（值 $1$），Right 横放剩单行 $1\times2$（值 $-1$），故
\[
 v(2\times2)=\{1\mid -1\}=\pm1 ,
\]
是一枚开关：谁先手谁赢。
\end{itemize}
把前几个值排成一行：$v(1\times1)=0$（谁都不能放），$v(1\times2)=-1$，$v(1\times3)=-1$，$v(1\times4)=-2$，$v(1\times5)=-2$，$v(1\times6)=-3$，恰是 $-\lfloor n/2\rfloor$。更大的棋盘（如 $3\times n$、$8\times8$ 整盘）值要靠计算机枚举，《Winning Ways》里有成套分析；本书只取最小情形演示方法。

\begin{problem}{蓝红伐木游戏（Blue--Red Hackenbush）}
棋盘改为一般图形：边接在地面或别的边上，Left 只砍蓝边、Right 只砍红边，被砍边的悬空部分整体倒下，无法行动者负。求给定图形的值。
\end{problem}

竹竿只是最简单的图形。《Winning Ways》的一个基础结论是：\textbf{任意蓝红伐木局面的值都是数}——蓝红世界里没有模糊性，这正是它适合给超现实数“接生”的原因。若加入双方都能砍的\textbf{绿边}，模糊性立刻回来：单独一条绿边 $\{0\mid 0\}=\ast$，又回到第 5 章的世界。蓝、红、绿三色版伐木游戏因此可以看成“数的世界 + 模糊游戏”的完整舞台，也是检验 6.4 节四种结果分类的最佳沙盘。

\paragraph{绿竿值多少.} 单条绿边就是 $\ast$。两节绿竿：Left 砍底，整竿倒下，Right 面对空局面直接负；Right 对称地先手也胜。其左右选项同为 $\{0,\ast\}$，记作 $\{0,\ast\mid 0,\ast\}$，与 $0$ 模糊。更有意思的是绿与数的对峙：$1+\ast$ 中 Left 先手把 $\ast$ 走成 $0$，Right 面对 $1$ 无事可做；Right 先手同样只能把 $\ast$ 走成 $0$，Left 再把 $1$ 走成 $0$。两种先后手都是 Left 胜——\textbf{任何数都压得过先手权}（“数与模糊游戏比较”的通例，见《Winning Ways》）。同理 $\tfrac14+\uparrow+\downarrow>0$ 是定理~\ref{thm:up} 微小性的直接应用。于是绿边的加入让伐木游戏变成“数的部分 + 模糊的部分”的混合体：数的部分按算术结算，模糊部分用第 5 章的工具处理——这正是三色伐木被视为“连接第 5、6 两章之桥”的原因。

\begin{pitfall}[开关不是数]
$\pm1=\{1\mid -1\}$ 不满足“左选项 $<$ 右选项”，不能参与“比大小”意义上的算术；比较热博弈时用“均值”只是启发式，不是相等关系（例~\ref{ex:hotstrat} 中均值 $\tfrac32>0$ 的局面 Right 先手照样赢）。若一道题只关心胜负而出现了开关，先把开关当成独立分支按先后手讨论，别急着套数轴。
\end{pitfall}

\section{拓展小结：一张值表}
把本章出现过的值集中成一张表，复习时先盖住后两列，自己说出每个游戏的值与结果类。

\begin{center}
\begin{tabular}{llll}
\toprule
游戏 & 值 & 结果类 & 一句话\\
\midrule
$\{\mid\}$（空局面） & $0$ & 零游戏，后手胜 & 无事可做\\
$\{0\mid\}$（L 一颗石子） & $1$ & 正数，L 必胜 & L 领先一步\\
$\{\mid 0\}$（R 一颗石子） & $-1$ & 负数，R 必胜 & R 领先一步\\
$\{0\mid 1\}$（蓝红竹） & $\tfrac12$ & 正数，L 必胜 & 半目优势\\
$\{0\mid \tfrac12\}$（蓝红红） & $\tfrac14$ & 正数，L 必胜 & 四分之一目\\
$\{1\mid -1\}$（Domineering $2\times2$） & $\pm1$ & 模糊，先手胜 & 开关，不是数\\
$\{0\mid 0\}$（单条绿边） & $\ast$ & 模糊，先手胜 & 先手权\\
$\{0\mid \ast\}$（up） & $\uparrow$ & 正，L 必胜 & 一丝真实优势\\
$\{\ast\mid 0\}$（down） & $\downarrow$ & 负，R 必胜 & $-\uparrow$\\
$\uparrow+\downarrow$ & $0$ & 零游戏，后手胜 & 镜像抵消\\
\bottomrule
\end{tabular}
\end{center}

读表三句话：值为\textbf{数}的局面是“冷”的，直接做加减法、比大小就能判定胜负；值为\textbf{模糊}的局面必须分“谁先手”讨论；\textbf{微小值}介于两者之间——比 $0$ 强、比任何正数弱，而且它本身不是数（$\uparrow$ 的右选项 $\ast$ 不是数，不满足数的构造规则）。这张表也是 6.4 节“四种结果”分类的完整对照：正数与负数对应 L、R 的“专属优势”，零与模糊对应先后手的攻防。

\section{分层例题与解析}
以下例题把本章拓展的三件工具——单纯性法则、开关、微小值——各练一遍；解析默认读者已读过本章相应小节。

\begin{problem}{例 1：混色竹竿求值}
求蓝红红、红蓝两根竹竿的值，并判定局面“蓝红红 $+$ 红蓝”（两根并玩）谁必胜。
\end{problem}

\begin{problem}{例 2：Domineering 小局拼合}
求 $1\times4$、$2\times1$ 与 $2\times2$ 的值（直接引用上一节的结论亦可），并判定两个局面谁必胜：(a) $2\times1+1\times4$；(b) $2\times2+1\times2$。
\end{problem}

\begin{problem}{例 3：验证 $\uparrow$ 的三条性质}
验证 $\uparrow>0$、$\uparrow\parallel\ast$、$\uparrow+\downarrow=0$。
\end{problem}

\begin{problem}{例 4：兑现支票游戏}
定义一列游戏 $X_0=0$，$X_{n}=\{X_{n-1}+1\mid X_{n-1}-1\}$。求 $X_1,X_2,X_3$ 的值，并判定局面 $X_2+\tfrac12$ 与 $0$ 的关系。
\end{problem}

\begin{problem}{例 5：竹竿拼合与镜像}
求局面“蓝 $+$ 红”（两根单色竹竿并玩）与局面“红蓝 $+$ 蓝红红”的值，并判定胜负。
\end{problem}

\begin{problem}{例 6：Domineering 拼合与均值的边界}
求 $2\times1+2\times1$ 与 $2\times2+2\times1$ 的值（各分量值见本章 Domineering 一节），并判定胜负。
\end{problem}

\begin{problem}{例 9：数与模糊的对峙}
判定四个局面与 $0$ 的关系：(a) $\ast+1$；(b) $\ast+\ast$；(c) $\tfrac14+\uparrow+\downarrow$；(d) 蓝红红 $+$ 蓝红 $+$ 红（三根竹竿并玩）。
\end{problem}

\hint 数的部分与模糊的部分分开结算；(b) 想想两个先手权；(d) 全是数，直接做算术。

\begin{problem}{例 10（自拟）：判断题组}
判断真伪并说明理由：
(1) $\uparrow+\uparrow+\ast>0$；
(2) $\tfrac12+\ast$ 与 $0$ 模糊；
(3) $\uparrow+\ast$ 与 $0$ 模糊；
(4) $\uparrow+\uparrow+\downarrow$ 的值为 $\uparrow$；
(5) 蓝红红 $+\ast$ 与 $0$ 模糊。
\end{problem}

\hint 每题先分清“数 / 模糊 / 微小”三种成分；(1)(3) 要分先后手逐线验证，别背结论。

\begin{problem}{例 11（自拟）：换色竹竿}
把三节竹竿自下而上接成“蓝、红、蓝”，求其值并判定结果类。
\end{problem}

\hint 逐个颜色考虑“砍这里会剩下什么”，再按左右选项的支配关系取极值。

\begin{problem}{例 7：开关的加法}
求 $\pm2+\pm1$ 的值（其中 $\pm k=\{k\mid -k\}$），并判定局面与 $0$ 的关系。
\end{problem}

\begin{problem}{例 8：结果类分类}
对下列游戏，指出它与 $0$ 的关系（大于、小于、等于、模糊），并说明哪些是数：$0$；$\tfrac14$；$\pm1$；$\ast$；$\uparrow$。
\end{problem}

\chapter{纳什均衡与零和博弈}
\section{纳什均衡}
静态非合作博弈问题的局部最优策略称为纳什均衡。

\begin{definition}[纳什均衡]
一个策略组合是\textbf{纳什均衡}，若其中任何一个玩家单方面改变策略，都不能提高他的收益（局部最优）。
\end{definition}

先解释“局部”二字：纳什均衡只要求“没有人想单方面反悔”，并不保证整体收益最大，甚至不保证均衡是唯一的；它刻画的是理性博弈中一个“稳定”的拍板状态——给定别人的选择，我已经做了我能做的最好的事。

“策略”指的是什么？最简单的想法是：每名玩家的策略就是固定一种决策（\textbf{纯策略}）。然而，按照这个定义，纳什均衡不一定存在。

\begin{example}[石头剪刀布没有纯策略均衡]
胜者收益 $+1$，败者收益 $-1$，平局双方收益均为 $0$。收益矩阵（行是 $A$ 的出拳，列是 $B$ 的出拳，格子里是 $A$ 的收益；$B$ 的收益恰为其相反数）为
\[
 W=\begin{pmatrix}0&1&-1\\-1&0&1\\1&-1&0\end{pmatrix}.
\]
任取非胜的一方，固定另一方的策略，这一方一定可以通过反悔取胜，所以不存在（纯策略）纳什均衡。换言之：在纯策略的世界里，石头剪刀布没有“稳定的拍板状态”，谁先亮拳谁吃亏。
\end{example}

根据生活经验，石头剪刀布博弈的最优策略是随机出拳，而不是一直出同样的拳。这就引出了一种更复杂的策略。

\begin{definition}[纯策略与混合策略]
固定一种决策的策略称为\textbf{纯策略}；为每种决策分配一个使用概率、每次依概率随机决策的策略称为\textbf{混合策略}。
\end{definition}

在混合策略下，纳什均衡比较的是玩家反悔前后的\textbf{期望}收益——决策随机化之后，“收益”不再是确定值，比较的基准自然变成期望。可以验证石头剪刀布中唯一的均衡是双方各以 $\tfrac13$ 的概率出三种拳：此时无论对方出什么，自己的期望收益都是 $0$，任何偏离（例如偏向石头）都会被对方用“专门出布”惩罚到负期望，因此没有人愿意单方面改变。

\begin{theorem}[纳什定理]\label{thm:nash}
混合策略纳什均衡一定存在。
\end{theorem}

纳什定理告诉我们：把策略空间放宽到概率分布后，“稳定拍板点”总是存在；这一定理是 1950 年纳什用角谷不动点定理证明的，也是他获得诺贝尔经济学奖的工作之一。

纯策略纳什均衡太简单了（枚举决策组合即可判定），下文以混合策略纳什均衡为主线；唯一需要两者并存的场合是双人零和博弈——那里纯策略均衡对应鞍点、混合均衡给出唯一的值，见后面“零和博弈”与“Minimax 定理”两小节。

纳什均衡有什么用？在 OI 中几乎没用。理由如下：
\begin{itemize}
\item 因为纳什均衡仅仅是一个局部最优策略，所以对于同一个博弈问题，可能存在多个毫不相关、无法横向对比的纳什均衡。题目若要求“求出最优策略”，需要的是全局最优；而纳什均衡连“哪个均衡更好”都说不清。
\item 更糟糕的是，求解纳什均衡的复杂度很高。\footnote{参考资料：\url{https://people.cs.pitt.edu/~kirk/CS1699Fall2014/lect4.pdf}。相关经典结论：判定是否存在多个纳什均衡是 NP-Complete 的（Gilboa--Zemel, 1989）；求任意一组纳什均衡是 PPAD-Complete 的（Chen--Deng, 2006）。}
\begin{itemize}
\item 判定是否存在两个纳什均衡是 NP-Complete 的：3-SAT 可以规约到双人匿名博弈的情况。
\item 求任意一组纳什均衡是 PPAD-Complete 的：求布劳威尔不动点可以规约到双人博弈的情况。
\end{itemize}
\end{itemize}
所以，在 OI 中可能有用的纳什均衡问题只有：双人零和博弈。下面先补一点零和博弈本身的性质，Minimax 定理紧随其后。

\section{零和博弈}
零和的意思是收益总和恒定——在全篇的约定里通常就是 $0$：一方所得即另一方所失，$B$ 的收益恰好是 $A$ 的收益的相反数。于是整场博弈只需要一张表就够了：$A$ 的收益矩阵 $W$，$B$ 的全部心思就是最小化它。注意 1.2 节把“非合作（零和）”列为组合游戏的特征，这里讨论的则是静态的两人零和博弈，视角不同，但“利益完全对立”的内核是同一个。

\begin{example}[猜硬币（Matching Pennies）]\label{ex:pennies}
两人同时各亮一枚硬币，同面则 $A$ 赢 $1$ 元、异面则 $A$ 输 $1$ 元，收益矩阵（A 视角）为
\[
 W=\begin{pmatrix}1&-1\\-1&1\end{pmatrix}.
\]
这个矩阵没有纯策略纳什均衡：$A$ 出正面时 $B$ 想出反面，$A$ 出反面时 $B$ 想出正面，任何“固定亮法”都会被人针对。所以双方只能随机：各以 $\tfrac12$ 的概率出正反，均衡值 $0$。矩阵是反对称的（$W_{i,j}=-W_{j,i}$，决策集相同），对称性保证了没人能占便宜、均衡值只能是 $0$；石头剪刀布（上一节开头那张 $3\times3$ 矩阵）就是它的三决策版本。
\end{example}

这里附带一个常用技巧：给收益矩阵整体加一个常数或乘一个正数，所有均衡策略都不变，只是博弈值随之平移或缩放——因此分析时可以把矩阵平移成非负、放大成整数，怎么方便怎么来（这也是下一小节线性规划写法能成立的伏笔）。

如果纯策略下就有“互相锁死”的格子，事情会简单得多。设 $A$ 能保底的收益是 $\max_i\min_j W_{i,j}$（A 先挑行，B 再挑最差列），$B$ 能把 $A$ 压到的上限是 $\min_j\max_i W_{i,j}$（B 先挑列，A 再挑最好行），显然恒有 $\max\min\le\min\max$。

\begin{definition}[鞍点]
若存在格子 $(i^*,j^*)$ 同时实现 $\max_i\min_j W_{i,j}$ 与 $\min_j\max_i W_{i,j}$，则称 $(i^*,j^*)$ 为\textbf{鞍点}。
\end{definition}

若鞍点存在，双方直接选 $i^*,j^*$ 就是纯策略纳什均衡。

\begin{example}[鞍点]
\[
 W=\begin{pmatrix}1&1\\2&0\end{pmatrix}
\]
中，$\max_i\min_j W_{i,j}=1$（A 选第一行，无论 B 怎么选都保底 $1$），$\min_j\max_i W_{i,j}=1$（B 选第二列，无论 A 怎么选都最多给 $1$），格子 $(1,2)$ 就是鞍点。鞍点的直观形象是“行最小、列最大”：A 不敢换行（换了可能更差），B 不敢换列（换了 A 可能更好），谁先偏离谁吃亏。反过来看猜硬币，$\max\min=-1$ 而 $\min\max=1$，空隙里全是“想反悔”的冲动，纯策略自然稳不住。
\end{example}

零和博弈最可爱的地方在于它有唯一的值：Minimax 定理保证 $\max\min=\min\max$，于是无论纯策略还是混合策略，所有纳什均衡都给出同一个期望收益 $v$，不存在一般博弈里“多个均衡、无从比较”的尴尬。这也正是前面说“OI 里有用的纳什均衡只剩双人零和”的底气：题目要的“最优策略下的期望收益”在零和博弈里是个良定义的数，换成非零和博弈，连“最优”都说不清。鞍点存在就直接用纯策略，没有鞍点（如猜硬币）就上混合策略——具体怎么解，看下面的 Minimax 定理与图解法。

\section{Minimax 定理}
\begin{problem}{Minimax 问题}
AB 两人博弈，A 有 $n$ 种决策，B 有 $m$ 种决策。如果 A 选择决策 $i$，B 选择决策 $j$，则游戏的分数是 $W_{i,j}$。$W$ 称为游戏的收益矩阵。A 要最大化期望分数，B 要最小化期望分数，求游戏的期望分数。
\end{problem}

\begin{theorem}[Minimax 定理]\label{thm:minimax}
这个问题的答案是唯一的。
\end{theorem}

先感受一下为什么“唯一”不平凡：B 想最小化分数，但他不知道 A 会怎么随机；A 想最大化分数，也不知道 B 怎么随机。两人都在跟“对方的未知策略”博弈。Minimax 定理说的是：这种相互猜测最终收敛到一个确定值——存在混合策略 $(p,q)$ 与常数 $v$，使得 A 无论怎么走，在 $q$ 下的期望至多 $v$；B 无论怎么走，在 $p$ 下的期望至少 $v$，而 $v$ 就是博弈的“公平价格”。注意这里 $v$ 同时是 A 的保底与 B 的天花板，两者相等是定理的全部内容。

\paragraph{A 的角度（下界）.} A 可以选择一个概率分布 $p_{1\ldots n}$，来最大化在 B 的任意策略下游戏的期望分数。具体地说，A 要最大化“最坏情况”：无论 B 选哪个纯策略 $j$，期望分数 $\sum_i p_iW_{i,j}$ 都必须 $\ge ans$。这个最坏情况下的最大值是答案的一个下界——A 至少能保证这么多。

\paragraph{B 的角度（上界）.} 对称地，B 可以选择一个概率分布 $q_{1\ldots m}$，来最小化在 A 的任意策略下游戏的期望分数：无论 A 选哪个纯策略 $i$，期望分数 $\sum_j q_jW_{i,j}$ 都必须 $\le ans'$。这是答案的一个上界——B 最多让 A 拿到这么多。

只要证明上界等于下界，就能证明原问题的答案唯一。若下界 $<$ 上界，两人各自的“保底”之间有空隙，说明存在相互占优的余地，均衡不存在；定理断言空隙不存在。

考虑 A 的角度。问题可以表示成线性规划的形式：设答案为 $ans$。
\[
 \begin{aligned}
 \max\quad & ans\\
 \text{s.t.}\quad & \sum_{i=1}^{n}p_i\le1\\
 & \forall j\in[1,m],\ ans-\sum_{i=1}^{n}p_iW_{i,j}\le0
 \end{aligned}
\]

\begin{pitfall}[$\sum p_i\le1$ 还是 $\sum p_i=1$]
严格地说，对一般的实收益矩阵，约束应写成 $\sum_i p_i=1$（概率质量必须恰好用完）；这里写 $\le1$ 相当于允许 A“弃权”部分概率质量，只有当收益矩阵非负（或先整体平移成非负）时两种写法才给出同一最优值。整体给矩阵加常数只平移博弈值、不改变均衡策略，乘正数同理——所以遇到一般矩阵，先平移成非负即可放心使用这套写法。OI 中的双人零和题（胜率、金钱收益）矩阵天然非负，更无此顾虑。
\end{pitfall}

考虑其对偶形式：
\[
 \begin{aligned}
 \min\quad & ans'\\
 \text{s.t.}\quad & \sum_{j=1}^{m}q_j\ge1\\
 & \forall i\in[1,n],\ ans'-\sum_{j=1}^{m}q_jW_{i,j}\ge0
 \end{aligned}
\]
观察到这就是 B 的答案对应的线性规划！证毕。这里的“证毕”其实调用了线性规划的\textbf{强对偶定理}：对偶问题的最优值等于原问题的最优值。而恰好 B 的“最小化最坏情况”写出来就是 A 的线性规划的对偶，两个最优值相等，上下界闭合，Minimax 定理得证。从理论层面看，Minimax 定理是“LP 强对偶”在博弈论里的一个化身。

使用求解线性规划的算法也可以构造出一组策略：解出上述 LP（或其任何等价形式）后，最优解 $p,q$ 本身就是双方的均衡混合策略。不过 OI 中一般用不到通用的 LP 求解器，实战的出路在下一小节。

\section{能不能不用线性规划？}
对于一般情况，不行。对于一般题目，可以。

实战中遇到的双人零和博弈，一般每人都只有两种决策。这种情况下可以使用\textbf{图解法}：设 A 使用第一种决策的概率为 $p$（则第二种决策的概率为 $1-p$）。观察到对于 B 的每种决策，当 B 只使用这种决策时，A 的期望分数是一个关于 $p$ 的一次函数：
\[
 f_j(p)=pW_{1,j}+(1-p)W_{2,j}.
\]
画出这些一次函数的图像，容易通过分类讨论求出纳什均衡。具体做法如下：
\begin{itemize}
\item 对每个 $j$ 画直线 $y=f_j(p)$（$p\in[0,1]$）。A 的保证收益是“最坏情况”，即所有直线中的\textbf{下包络} $\min_j f_j(p)$——B 一定会挑最有利于他的那条线来对付 A。
\item 纳什均衡中的 $p$ 就是下包络的最高点 $p^*=\arg\max_{p\in[0,1]}\min_j f_j(p)$，均衡值 $v=\min_j f_j(p^*)$。
\item 最高点只可能出现在两种地方：两条直线的交点，或区间端点 $p=0,1$。若 $f_1,f_2$ 的斜率符号相反（B 的两种决策对 A 的威胁一增一减），下包络呈“倒 V”形或“V”形，最大值点自然在两线交点处，解方程 $f_1(p)=f_2(p)$ 即可；否则退化到端点，检查哪个端点收益更大。
\end{itemize}

\begin{center}
\begin{tikzpicture}[scale=3.2,>=Stealth]
\draw[->] (-0.05,0) -- (1.15,0) node[right] {$p$};
\draw[->] (0,-0.05) -- (0,1.0);
\draw (1,-0.02) node[below] {$1$} -- (1,0.02);
\draw (0.015,1) node[left] {$y$} -- (-0.015,1);
\draw[thick,ink!70] (0,0.9) -- (1,0.2) node[right] {\small $f_1(p)$};
\draw[thick,accent] (0,0.3) -- (1,0.8) node[right] {\small $f_2(p)$};
\draw[very thick,warn] (0,0.3) -- (0.7,0.55) -- (1,0.2);
\filldraw[warn] (0.7,0.55) circle (0.6pt) node[above=2pt] {\small $(p^*,v)$};
\end{tikzpicture}\\[2pt]
{\small 图解法：实心折线是两条期望直线的下包络，其最高点 $(p^*,v)$ 即 A 的均衡混合概率与博弈值。}
\end{center}

当交点 $p^*\in(0,1)$ 时，B 的均衡混合策略 $q$ 也可以用对称方法求出，或者用“在均衡点处 B 两种纯决策的期望相等”这一条件反解。例题 [CF98E] Help Shrek and Donkey 和 [THUPC 2023 Pre] 欺诈游戏的解法都会落到这个流程上。

如果遇到了决策数很多的双人博弈，可以先想想能不能转化成每人只有两种决策的情况。这一步往往才是题目的真正难点：例如 [THUPC 2023 Pre] 欺诈游戏看似每方有 $n$ 种决策，其实可以用递推把“是否选择 $n$”提炼成 $2$ 种决策的博弈。

如果想不到，可以考虑使用一个网络上广为流传的乱搞做法：不妨设 $n\le m$。设 $w_j$ 表示如果 B 只使用决策 $j$，A 的期望分数。则如果 $w$ 不两两相等，那么直觉上可以“移多补少”，调整得到一个更优的方案。考虑直接列方程高斯消元：根据 $w$ 两两相等可以列出 $m-1$ 个方程，此外我们还有 $\sum p_i=1$，所以方程数量一定足够。为什么均衡时 $w$ 要两两相等？因为在均衡点处 B 对他用到的每个纯策略应当“无所谓”（期望相同），否则 B 会放弃收益差的那个策略；这个条件叫\textbf{无差异原理}，是手动解小型博弈的利器。因为题目往往有特殊性质，所以这样做很可能是对的。

根据 Itst 的说法，当 $n=m$ 时，如果从 A 的角度出发，方程的解满足 $p_i\ge0$，且从 B 的角度出发，方程的解满足 $q_i\ge0$，那解方程就是正确的。我们不会证明，但是这个判定条件似乎对解题并无帮助。

\section{图解法完整算例}
抽象流程说完了，本节把一个具体的 $2\times2$ 博弈从头算到尾，作为图解法的标准作业流程。\footnote{延伸资料：洛谷专栏《浅谈纳什均衡》（\url{https://www.luogu.me/article/pjhj07pz}），其中强调的“每个参与人的混合策略都使其余参与人任何纯策略的期望收益相等”正是本节反复使用的无差异原理；文中关于匿名博弈纳什均衡近似求解的讨论也值得一读。}

\begin{problem}{完整算例}
双方各两种决策，收益矩阵（A 视角）为
\[
 W=\begin{pmatrix}2&-1\\-1&1\end{pmatrix},
\]
求均衡混合策略与博弈值 $v$。
\end{problem}

\sol \paragraph{第一步：排除纯策略均衡（找鞍点）.} 行最小值为 $-1,-1$，故 $\max\min=-1$；列最大值为 $2,1$，故 $\min\max=1$。两者不等，无鞍点，必须混合。

\paragraph{第二步：写两条期望直线.} 设 A 以概率 $p$ 选第一行，则 B 各纯策略下 A 的期望为
\[
 f_1(p)=2p-(1-p)=3p-1,\qquad f_2(p)=-p+(1-p)=1-2p .
\]
\paragraph{第三步：找下包络的最高点.} $f_1$ 递增、$f_2$ 递减，包络呈“倒 V”形，交点满足 $3p-1=1-2p$，解得
\[
 p^{*}=\frac25,\qquad v=f_1(p^{*})=\frac65-\frac15=\frac15 .
\]
A 以 $\tfrac25$ 混合时，无论 B 出哪一列，A 的期望都是 $\tfrac15$——这就是 A 的保底。

\paragraph{第四步：反解 B 的混合（无差异原理）.} 设 B 以概率 $q$ 选第一列。均衡点处 A 对两行应当无所谓：
\[
 2q-(1-q)=-q+(1-q)\ \Longrightarrow\ 3q-1=1-2q\ \Longrightarrow\ q^{*}=\frac25,
\]
此时 B 把 A 恰好压到 $v=\tfrac15$。本例中 $p^{*}=q^{*}$ 纯属巧合（矩阵接近对称），一般并不相等。

\paragraph{第五步：验算.} 若 A 偏离到 $p=0$，B 用第一列把 A 压到 $-1<v$；若 B 偏离到全用第二列，A 拿到 $1>v$。双方都无利可图，均衡成立。整个流程——排除鞍点、画包络、解交点、无差异反解、验算——对任何 $2\times n$ 图解法题目都是同一套。

\begin{pitfall}[端点均衡：当交点不在开区间内时]
下包络的最高点未必是两条直线的交点：若两线同向（B 的两种策略对 A 的威胁同增同减），最高点落在 $p=0$ 或 $p=1$，均衡退化为纯策略——此时通常存在被支配的纯策略，先删支配再解往往更省事。判据：交点 $p^{*}\notin[0,1]$ 或两线斜率同号时，直接比较两个端点的包络值。
\end{pitfall}

\section{拓展：冯·诺依曼 1928——Minimax、对偶与扑克}
1928 年，二十出头的冯·诺依曼发表《Zur Theorie der Gesellschaftsspiele》（论博弈论）\footnote{Mathematische Annalen 第 100 卷，295--320 页；1926 年 12 月在格丁根数学学会首报，\url{https://doi.org/10.1007/BF01448847}。}，第一次对一般二人零和博弈证明了定理~\ref{thm:minimax}。1953 年他在给弗雷歇的公开答复里写道：“就我所见，没有这条定理，就没有现代意义上的博弈论……在 Minimax 定理被证明之前，我以为没有什么值得发表。”（ Econometrica 上关于波雷尔批注的通信。）事实上波雷尔从 1921 年起就在一系列短文里提出混合策略的想法，并在极小的例子上猜测这条定理，但始终没能证明；弗雷歇 1953 年为波雷尔争优先权，冯·诺依曼以此作答。7.3 节的证明（线性规划对偶）在历史上出现得晚得多：1928 年的原始证明纯组合且相当繁难，后来经不动点方法（冯·诺依曼 1937、纳什 1950）与线性规划对偶（丹齐格 1951）被反复重证、日益透明。据丹齐格回忆，1947 年他向冯·诺依曼介绍线性规划时，对方听完当即回敬了一小时即兴讲演，讲的正是对偶性以及它与矩阵博弈的等价——7.4 节讨论的问题，源头就是这次会面。

论文前后，冯·诺依曼用刚诞生的 Minimax 理论做了一件至今仍令人叫绝的事：分析扑克。他把“诈唬”从心理花招变成了数学必然——Kuhn 与 Tucker 后来称之为“连外行都能欣赏的突破”。本节用本章的工具把这个模型完整解掉。

\paragraph{模型.} 规则：双方各下底注 $1$ 元；各从 $[0,1]$ 均匀抽一个隐藏数（A 拿 $x$，B 拿 $y$），数大者胜，各自只见自己的数；A 先行动：弃牌（输掉底注，收益 $-1$）或下注 $a$ 元；若 A 下注，B 或弃牌（A 赢得 B 的底注 $+1$），或跟注（补足 $a$ 元后摊牌，大者全收）。不许加注。这是零和博弈，但信息不完全——谁也不知道对方的数。

\paragraph{均衡.} 设 B 用阈值策略：$y\ge c$ 才跟注。A 拿 $x<c$ 的牌下注：B 以概率 $c$ 弃牌送钱、以概率 $1-c$ 跟注且必胜，收益为 $c-(1-c)(1+a)=c(2+a)-(1+a)$；弃牌收益 $-1$。两者相等的临界值即
\[
 c=\frac{a}{2+a}.
\]
$x\ge c$ 时下注收益为 $c+(1+a)(2x-c-1)$，随 $x$ 递增，故 A 全部价值下注；$x<c$ 的手牌在均衡处诈唬与弃牌无差异，下注频率由 B 的无差异条件定：B 拿 $y=c$ 的牌，摊牌胜率是 A 诈唬牌占下注范围的份额 $\tfrac{pc}{pc+1-c}$（$p$ 为诈唬频率），令其等于 $\tfrac{a}{2(1+a)}$，解出
\[
 p=\frac{2}{2+a}.
\]
三类收益加总——弃牌区每手 $-1$、诈唬区每手恰为 $-1$、价值下注区积分恰为 $c(1-c)$——博弈值为
\[
 v=-c+c(1-c)=-c^{2}=-\Bigl(\frac{a}{2+a}\Bigr)^{2},
\]
对 A 为负：被迫先行动、先亮出“我下注了”这个信息的一方吃亏。

\begin{center}
\begin{tabular}{cccc}
\toprule
赌注 $a$ & B 跟注阈值 $c$ & A 诈唬频率 $p$ & 博弈值 $v$\\
\midrule
$\tfrac12$ & $\tfrac15$ & $\tfrac45$ & $-\tfrac1{25}$\\
$1$ & $\tfrac13$ & $\tfrac23$ & $-\tfrac19$\\
$2$ & $\tfrac12$ & $\tfrac12$ & $-\tfrac14$\\
\bottomrule
\end{tabular}
\end{center}

\paragraph{为什么必须诈唬.} 假如 A 从不诈唬，只在好牌下注，B 的最优回应是“牌不够大一律弃”，A 的价值下注便一无所获；于是 A 必须用一部分垃圾牌冒充强牌，把 B 逼回赌桌，而诈唬频率恰好被调到 B 愿意跟注为止——7.5 节的无差异原理在这里两头各捏一刀：A 的阈值由自己的无差异定，A 的频率由 B 的无差异定。注意赌注越大，B 越谨慎（$c$ 升高）、A 诈唬越罕见（$p$ 下降）但区间越宽、先手劣势 $c^{2}$ 也越大。

\begin{pitfall}[诈唬不是越多越好]
学生常把“均衡要诈唬”读成“诈唬越多越好”。练习 8 给出反例：$a=1$ 时把频率提到 $1$，B 的阈值降到 $\tfrac14$，每手垃圾牌比弃牌多亏 $\tfrac14$ 元。反向误读同样要防——“均衡里有诈唬，所以我随时可以诈”：离开均衡频率，诈唬立刻从必需品变成送给对手的礼物。
\end{pitfall}

这对策略在 7.1 节的意义下恰是纳什均衡：单方面多诈、少诈、放宽或收紧阈值都只会更差。冯·诺依曼比纳什早二十二年，就在零和世界里把“各玩各的最优回应”走到了尽头；2017 年击败职业选手的扑克 AI Libratus，在算力层面延续的仍是“按均衡频率诈唬”这条 1928 年的思路。

\section{分层例题与解析}

以下例题的共性：真实博弈被压缩成一个 $2\times2$ 零和子博弈，再用本章的图解法求解；博弈论含量集中在“找到正确的 $2\times2$ 子博弈”这一层。


\begin{problem}{练习 1（自拟）：三线包络}
双方博弈，A 有两种决策、B 有三种决策，收益矩阵（A 视角）为
\[
 W=\begin{pmatrix}2&-1&1\\-1&1&1\end{pmatrix},
\]
求 A 的均衡混合策略、博弈值，并说明 B 的第三种决策在均衡中的地位。
\end{problem}

\hint 三条期望直线逐对求交点，找出真正构成下包络的两条；整体位于包络上方的直线对应的纯策略会被闲置。

\begin{problem}{练习 2（自拟）：非对称石头剪刀布}
收益规则：平局 $0$；A 赢常规一局得 $1$、输一局得 $-1$；唯独“石头砸剪刀”这一手 A 赢得 $2$（输时仍为 $-2$）。收益矩阵（行 $A\in$\{石头,布,剪刀\}，列 $B$ 同序）为
\[
 W=\begin{pmatrix}0&-1&2\\1&0&-1\\-2&1&0\end{pmatrix},
\]
求双方均衡混合策略与博弈值。
\end{problem}

\hint 矩阵反对称（$W_{i,j}=-W_{j,i}$），均衡值必为 $0$；用无差异原理列两个方程再加归一化即可。

\begin{problem}{练习 3（自拟）：不对称猜硬币}
两人同时亮硬币，同面 A 赢 $1$ 元、异面 A 输 $1$ 元；但若 A 出反面而 B 出正面，A 输 $2$ 元。收益矩阵（行 \{正,反\}，列 \{正,反\}）为
\[
 W=\begin{pmatrix}1&-1\\-2&2\end{pmatrix},
\]
求均衡与博弈值。
\end{problem}

\hint 先查鞍点；再按“下包络最高点”解 $p$，用无差异原理反解 $q$。

\begin{problem}{练习 4：一般 $2\times2$ 混合均衡的闭式解}
设收益矩阵 $W=\begin{pmatrix}a&b\\ c&d\end{pmatrix}$ 且无鞍点，推导均衡混合概率与值的一般公式，并用它复算完整算例（$a=2,b=-1,c=-1,d=1$）与练习 3。
\end{problem}

\hint 对 B 的两列写无差异方程解 $p^{*}$；对 A 的两行写无差异方程解 $q^{*}$；代回任一行求 $v$。

\begin{problem}{练习 5（自拟）：鞍点检索}
收益矩阵（A 视角）为
\[
 W=\begin{pmatrix}3&2&4\\1&2&0\\5&2&1\end{pmatrix},
\]
判断是否存在纯策略纳什均衡（鞍点）；若有，给出双方的选择与博弈值。
\end{problem}

\hint 分别算 $\max_i\min_j W_{i,j}$ 与 $\min_j\max_i W_{i,j}$，相等即有鞍点。

\begin{problem}{练习 6（自拟）：端点均衡与鞍点}
收益矩阵（A 视角）为
\[
 W=\begin{pmatrix}1&-2\\1&1\end{pmatrix},
\]
用图解法求均衡，并解释“最高点落在端点”意味着什么。
\end{problem}

\hint 写两条期望直线，找下包络最高点；端点均衡往往对应鞍点，回头用 $\max\min$ 与 $\min\max$ 验证。

\begin{problem}{练习 7（自拟）：小赌注的扑克}
取 $a=\tfrac12$，求 B 的跟注阈值 $c$、A 的诈唬频率 $p$ 与博弈值 $v$，并与 $a=1$ 比较：赌注变小后，B 更愿意跟注还是更谨慎？A 的诈唬更频繁还是更罕见？
\end{problem}

\hint 代入 $c=\tfrac{a}{2+a}$、$p=\tfrac{2}{2+a}$、$v=-\bigl(\tfrac{a}{2+a}\bigr)^{2}$。

\begin{problem}{练习 8（自拟）：诈唬过头的代价}
设 $a=1$，A 改为“所有手牌都下注”。求 B 的最优跟注阈值，并算出 A 的垃圾牌下注相对弃牌的净亏损，说明均衡频率 $\tfrac23$ 为什么不能再高。
\end{problem}

\hint 面对“全下注”，B 在哪条 $y$ 上跟注与弃牌无差异？

\begin{problem}{[CF98E] Help Shrek and Donkey}
史瑞克和驴决定玩一种叫做 YAGame 的牌类游戏。规则非常简单：最开始，史瑞克手中有 $m$ 张牌，驴手中有 $n$ 张牌（两位玩家都看不到对方手中的牌），还有一张牌正面朝下地放在桌子上，两位玩家同样也看不到。因此，游戏开始时共有 $m+n+1$ 张牌。此外，玩家们知道整副牌由哪些牌组成，也清楚自己手中有哪些牌（但他们不知道哪一张牌在桌子上，以及对方手中有哪些牌）。游戏由史瑞克先手，接下来两位玩家轮流行动。每轮中，某位玩家可以进行以下两种操作之一：
\begin{itemize}
\item 尝试猜测桌子上的那张牌。如果他猜对了，游戏立即结束，他获胜。如果猜错了，游戏同样立即结束，但对方获胜。
\item 指定整副牌中的任意一张牌。如果对方手中有这张牌，则必须亮出来并将其移出本局（这张牌不再参与后续游戏）。如果对方没有该牌，则直接说明。
\end{itemize}
如果两位玩家均采取最优策略，求史瑞克获胜的概率。

数据范围：$n,m\le1000$。
\end{problem}

\hint 设 $f_{n,m}$ 为 A（先手，$n$ 张牌）对 B（$m$ 张牌）的胜率，递归会互换攻守，需要“$1-f$”补集写法；决策二的分支里藏着一个 $2\times2$ 静态子博弈，用图解法在两直线交点处解均衡。

\cost 时间复杂度 $O(n^{2})$（$n,m\le1000$）。本题是“图解法 + 概率递推”的教科书组合：牌类游戏的不确定性被压缩成一个 $2\times2$ 收益矩阵，博弈论含量在“找到正确的 $2\times2$ 子博弈”这一层。

\src{https://codeforces.com/problemset/problem/98/E}

\begin{problem}{[THUPC 2023 Pre] 欺诈游戏}
在《LIAR GAME》中，小 E 看到了一个有趣的游戏。

这个游戏名叫《走私游戏》。游戏规则大概是这样的：一名玩家扮演走私者，一名玩家扮演检察官。走私者可以将 $x$ 日元（$x$ 为 $[0,n]$ 内的整数，由走私者决定）秘密放入箱子中，而检察官需要猜测箱子中的金额。假设检察官猜了 $y$（$y$ 也必须是整数）。如果 $x=y$，则走私失败，走私者一分钱也拿不到。如果 $x>y$，则走私成功，走私者可以从检察官那里拿走 $x$ 日元。如果 $x<y$，则走私失败，但是由于冤枉检察官需要赔付给走私者 $y/2$ 日元。游戏分有限回合进行，双方轮流做走私者和检察官。

可以证明，最优情况下每个回合走私者会采用同一种策略，检察官也会采用同一种策略。

小 E 想知道在一个回合中，双方的最优策略分别是什么，构造方案。答案对 $998244353$ 取模。

数据范围：$1\le n\le4\times10^{5}$。
\end{problem}

\hint 看似每人 $n$ 种决策，其实可以分步决策：先把“最大金额 $n$”的攻防单独拎出来，剩下的自动退化为 $n-1$ 的子问题；每层解一个 $2\times2$ 零和博弈。

\cost 时间复杂度 $O(n)$。本题演示了“压缩决策”这一招：$n$ 种决策的博弈若能找到某种“分层的最大元素”结构，就可以层层剥开化成 $2\times2$ 子博弈的递归。

\src{https://www.luogu.com.cn/problem/P9142}

\appendix
\chapter{核心判定模板与接口}
\section{模板使用说明}
本附录摘录全书反复出现的判定框架，完整 C++17 实现见资料包 \texttt{code/selftest.cpp}；摘录是教学模板，不是可直接提交的完整程序。所有模板共用一个前提：游戏在有限步内结束、无法行动者负（三分类模板额外支持平局）。自测程序对有限范围的随机与穷举数据，把每个公式与暴力搜索逐点比对。

\begin{lstlisting}[language=C++]
// mex：集合中未出现的最小非负整数。
int mex(const std::vector<int>& s) {
    std::vector<char> seen(s.size() + 1, 0);
    for (int x : s)
        if (0 <= x && x < (int)seen.size()) seen[x] = 1;
    int m = 0;
    while (seen[m]) ++m;
    return m;
}
\end{lstlisting}

\begin{lstlisting}[language=C++]
// 模板一：DAG 上的记忆化 SG（隐式转移版）。
// 约定：move(u) 枚举 u 的全部后继；图无环（游戏有限步结束）。
// g(u) == 0 恰为 P 态（定理 5.x）。
struct SGTable {
    std::function<void(int, const std::function<void(int)>&)> move;
    std::vector<int> sg;
    std::vector<char> done;
    SGTable(int n, decltype(move) m) : move(m), sg(n, 0), done(n, 0) {}
    int dfs(int u) {
        if (done[u]) return sg[u];
        done[u] = 1;
        std::vector<int> nxt;
        move(u, [&](int v) { nxt.push_back(dfs(v)); });
        return sg[u] = mex(nxt);
    }
};
\end{lstlisting}

\begin{lstlisting}[language=C++]
// 模板二：一般状态图（允许有环）上的三分类。
// 返回 col：0 = P 态，1 = N 态，2 = 平局（过河卒框架）。
// u 是 N 态 <=> 存在 P 态后继；u 是 P 态 <=> 全部后继都是 N 态。
std::vector<int> solve_game(int n,
                            const std::vector<std::pair<int,int>>& edges) {
    std::vector<std::vector<int>> g(n), rg(n);
    for (auto [u, v] : edges) { g[u].push_back(v); rg[v].push_back(u); }
    std::vector<int> col(n, 2), und(n);
    std::vector<char> done(n, 0);
    std::queue<int> q;
    for (int u = 0; u < n; ++u) {
        und[u] = (int)g[u].size();
        if (und[u] == 0) { col[u] = 0; done[u] = 1; q.push(u); }
    }
    while (!q.empty()) {
        int u = q.front(); q.pop();
        for (int p : rg[u]) {
            if (done[p]) continue;
            if (col[u] == 0) { col[p] = 1; done[p] = 1; q.push(p); }
            else if (--und[p] == 0) { col[p] = 0; done[p] = 1; q.push(p); }
        }
    }
    return col;  // 仍为 2 的点：双方都不离开未定子图 => 平局。
}
\end{lstlisting}

\begin{lstlisting}[language=C++]
// 模板三：阶梯 Nim 判定（定理 2.x）。
// 输入 a[1..n]：第 i 级台阶上的石子数；a[0] 弃用（地面）。
// 每次把第 i 级的至少一颗石子移到第 i-1 级，不能行动者负。
// 先手必胜 <=> 奇数级异或和不为 0（只有奇数级参与！）
bool stair_nim_wins(const std::vector<long long>& a) {
    long long s = 0;
    for (int i = 1; i < (int)a.size(); i += 2) s ^= a[i];
    return s != 0;
}
\end{lstlisting}

Nim、反 Nim、Nim-k 与反 Nim-k 的判定公式各只有几行，不再单独摘录：它们在自测程序中以函数形式出现，并直接与暴力搜索对拍。

\section{怎样复现验证}
在资料目录中运行：
\begin{lstlisting}[language=bash]
g++ -std=c++17 -O2 -Wall -Wextra code/selftest.cpp -o gt_selftest
./gt_selftest
\end{lstlisting}
自测程序完成五组有限范围验证：(1) “每次取 $1$ 或 $2$ 颗”的 SG 值周期 $3$；(2) Nim、反 Nim、Nim-k（$k=1,2$）与反 Nim-k（$k=1,2$）的判定公式对全部 $4$ 堆以内、每堆不超过 $5$ 颗的局面与暴力博弈搜索一致；(3) 三分类模板与不动点迭代在随机带环图上一致；(4) 阶梯 Nim 判定与暴力搜索一致；(5) 记忆化 SG 与显式转移表一致。暴力验证依赖的是“状态数小到可以穷举”这一事实，与正文证明互为印证而非替代；更大的参数范围仍以正文证明为准。

\end{document}
